\documentclass[12pt,a4paper]{article}

\usepackage{amsmath,amssymb,amsthm}
\usepackage{mathtools}
\usepackage{booktabs}
\usepackage{array}
\usepackage{hyperref}
\usepackage[margin=1in]{geometry}
\usepackage{enumitem}
\usepackage{graphicx}
\usepackage{xcolor}

\newtheorem{theorem}{Theorem}
\newtheorem{proposition}[theorem]{Proposition}
\newtheorem{lemma}[theorem]{Lemma}
\newtheorem{corollary}[theorem]{Corollary}
\newtheorem{conjecture}[theorem]{Conjecture}
\newtheorem{definition}[theorem]{Definition}
\newtheorem{observation}[theorem]{Observation}
\newtheorem{remark}[theorem]{Remark}

\newcommand{\R}{\mathbb{R}}
\newcommand{\C}{\mathbb{C}}
\newcommand{\Z}{\mathbb{Z}}
\newcommand{\Q}{\mathbb{Q}}
\newcommand{\N}{\mathbb{N}}
\newcommand{\ket}[1]{|#1\rangle}
\newcommand{\bra}[1]{\langle#1|}
\newcommand{\braket}[2]{\langle#1|#2\rangle}
\newcommand{\KS}{\mathrm{KS}}
\newcommand{\CK}{\mathrm{CK}\text{-}31}

\title{The Algebraic Landscape of Kochen--Specker Sets\\in Dimension Three}

\author{Michael Kernaghan}

\date{\today}

\begin{document}

\maketitle

\begin{abstract}
We present a computational survey showing that Kochen--Specker (KS) uncolorability in three-dimensional Hilbert space is controlled by the cancellation identities available in the coordinate alphabet. In all tested two-symbol alphabets $\mathcal{A} = \{0, \pm 1, \pm x\}$, KS sets arise only when $x$ satisfies one of two mechanisms: \emph{norm-2 cancellation} (an algebraic identity producing the value~2 from products of ring elements, as in $1+1=2$, $(\sqrt{2})^2=2$, or $\alpha\bar\alpha=2$) or \emph{phase cancellation} (a root-of-unity sum vanishing exactly, as in $1+\omega+\omega^2=0$). Alphabets where the smallest available cancellation involves terms of modulus $\geq 3$ produce orthogonal triples but not KS-uncolorability. This two-mechanism constraint explains why constructions cluster into six discrete algebraic islands among the fields tested. Three yield new KS configurations: the Heegner-7 ring $\Z[(1+\sqrt{-7})/2]$ (43 vectors, a genuinely new orthogonality graph), the golden ratio field $\Q(\varphi)$ (52 vectors, revealed only by cross-product completion), and $\Z[\sqrt{-2}]$ (a new algebraic realization of the Peres-type structure). Using SAT-based bipartite KS-uncolorability, we verify and extend the input counts of Trandafir and Cabello for bipartite perfect quantum strategies across all six islands. Cabello--Severini--Winter analysis shows that the algebraic substrate also determines the strength of quantum contextual advantage.
\end{abstract}

\section{Introduction}

The Kochen--Specker (KS) theorem~\cite{KochenSpecker1967} is one of the foundational no-go results of quantum mechanics, demonstrating that quantum observables cannot be assigned pre-existing definite values independently of the measurement context (for a comprehensive review, see Ref.~\cite{Budroni2022}). The theorem is typically proved by exhibiting a finite set of rays (one-dimensional subspaces) in a Hilbert space~$\mathcal{H}$ such that no consistent $\{0,1\}$-valued assignment exists satisfying the constraints imposed by~orthogonality.

In dimension $d = 3$, the smallest known KS set is the Conway--Kochen 31-vector set ($\CK$), reported by Peres~\cite{Peres1993} with coordinates drawn from the integer alphabet $\{0, \pm 1, \pm 2\}$. Despite decades of effort, no smaller construction has been found. The best known lower bound is 24 vectors, established independently by Li, Bright, and Ganesh~\cite{LiBrightGanesh2024} and by Kirchweger, Peitl, and Szeider~\cite{KirchwegarPeitlSzeider2023} using SAT-based methods with algebraic realizability checking (improving the earlier bound of 22 by Uijlen and Westerbaan~\cite{UijlenWesterbaan2016}), leaving a gap of 7 vectors that remains one of the central open problems in KS theory. Trandafir and Cabello~\cite{TrandafirCabello2025} recently proved that $\CK$ is rigid (unique up to unitary transformation in $\C^3$) and conjectured that 31 is optimal.

Previous computational approaches have primarily searched within fixed coordinate alphabets~\cite{Peres1991, Pavicic2004} or generated random abstract hypergraphs and tested realizability~\cite{LiBrightGanesh2024}. The algebraic structure of the coordinate systems underlying KS sets---specifically, which number fields can support them---has received less systematic attention, though Cortez, Morales, and Reyes~\cite{CortezMoralesReyes2022} proved that $\Z[1/6]$ is the minimal ring extension of $\Z$ exhibiting KS contextuality in dimension 3.

In this paper, we ask a structural question: \emph{what algebraic property of a coordinate alphabet controls whether it can support KS sets in dimension 3?} Through a systematic computational survey of quadratic fields $\Q(\sqrt{d})$ for $d = 2, \ldots, 30$, all nine imaginary quadratic fields with class number 1, cyclotomic fields, and golden ratio extensions, we find a sharp empirical answer: KS-uncolorability requires the alphabet to support a \emph{cancellation identity}---an algebraic relation among its elements that makes dot products vanish exactly---in which the largest participating term has modulus at most~$\sqrt{2}$ (norm-2 cancellation) or in which phases conspire to cancel (phase cancellation, as in roots of unity). Among all tested alphabets, these are the only two mechanisms that produce the dense triad networks required for KS-uncolorability, explaining why constructions cluster into six discrete ``algebraic islands'' among the fields tested.

The two-mechanism constraint unifies all known constructions: the integer identity $1 + 1 = 2$, the Peres identity $(\sqrt{2})^2 = 2$, the Eisenstein identity $1 + \omega + \omega^2 = 0$, and the Heegner-7 identity $\alpha\bar\alpha = 2$ are all instances of norm-2 or phase cancellation---the coordinate alphabet supports an algebraic relation whose terms have modulus at most $\sqrt{2}$, or whose phases sum to zero. Alphabets where the smallest available cancellation involves terms of modulus $\geq \sqrt{3}$ produce orthogonal pairs and triples but never enough interlocking triads for KS-uncolorability. Our main contributions are:

\begin{enumerate}[label=(\roman*)]
    \item An empirical classification of the cancellation identities that support KS-un\-color\-ability in two-symbol alphabets: all tested cases reduce to norm-2 or phase cancellation, with no intermediate mechanisms (Sections~\ref{sec:quadratic}--\ref{sec:heegner}).
    \item Two new KS configurations not previously reported: a 43-vector set from the Heegner-7 ring $\Z[\alpha]$ with $\alpha = (1{+}\sqrt{-7})/2$ (a genuinely new orthogonality graph), and a 52-vector set from the golden ratio field $\Q(\varphi)$, revealed only by cross-product completion (Sections~\ref{sec:heegner},~\ref{sec:golden}; Appendix~\ref{app:rays}).
    \item Evidence that the algebraic substrate has operational consequences: different islands yield different bipartite perfect quantum strategies (Section~\ref{sec:cabello}) and different Cabello--Severini--Winter contextual advantages (Section~\ref{sec:csw}), with no single island dominating on all measures.
    \item Supporting results: an exhaustive OCUS proof that no $\leq 30$-ray KS subset exists within the full 49-ray integer pool (Section~\ref{sec:exhaustive}), universal merge saturation across all six islands (Observation~\ref{obs:merge}), the pattern $6 \mid n$ for cyclotomic uncolorability (Section~\ref{sec:roots}), the $N(S)$ invariant connecting our islands to the minimal-ring framework of Cortez, Morales, and Reyes~\cite{CortezMoralesReyes2022}, and strong graph-invariant evidence for isomorphism across the norm-2 family (Proposition~\ref{prop:isomorphism}).
\end{enumerate}

\section{Preliminaries}\label{sec:prelim}

\subsection{Kochen--Specker sets and colorings}

\begin{definition}
A \emph{ray} in $\mathcal{H} = \C^d$ (or $\R^d$) is a one-dimensional subspace, represented by a nonzero vector up to scalar multiples. Two rays $v, w$ are \emph{orthogonal} if $\braket{v}{w} = 0$. An \emph{orthogonal triad} in $d = 3$ is a set of three mutually orthogonal rays forming a complete basis.
\end{definition}

\begin{definition}
A \emph{KS coloring} of a set of rays $S$ is a function $f: S \to \{0, 1\}$ (interpreted as ``red'' and ``green'') satisfying:
\begin{enumerate}[label=(\alph*)]
    \item \emph{Orthogonal exclusion}: if $v \perp w$ and $f(v) = 1$, then $f(w) = 0$;
    \item \emph{Triad completeness}: for every orthogonal triad $\{v_1, v_2, v_3\} \subset S$, exactly one $v_i$ has $f(v_i) = 1$.
\end{enumerate}
A set $S$ is a \emph{Kochen--Specker set} (or \emph{KS-uncolorable}) if no KS coloring exists.
\end{definition}

\begin{definition}
The \emph{orthogonality graph} $G(S)$ of a ray set $S$ has vertices corresponding to rays and edges connecting orthogonal pairs. A \emph{triad} corresponds to a triangle (3-clique) in $G(S)$. We use ``orthogonality graph'' when referring to the pair structure and ``orthogonality hypergraph'' when emphasizing the triad (3-uniform hyperedge) structure; both encode the same underlying ray configuration.
\end{definition}

\subsection{Coordinate alphabets and number fields}

A \emph{coordinate alphabet} $\mathcal{A} \subset \R$ (or $\C$) is a finite set of values from which ray coordinates are drawn.  (Pavicic et al.~\cite{Pavicic2019} call these ``vector components''; we adopt ``alphabet'' to emphasize the combinatorial generating role---one selects symbols from $\mathcal{A}$ and assembles them into three-coordinate ``words'' that represent rays.)  Given an alphabet $\mathcal{A}$, the induced ray set is
\[
S(\mathcal{A}) = \{ [v] : v \in \mathcal{A}^3 \setminus \{0\} \}
\]
where $[v]$ denotes the equivalence class of $v$ under scalar multiplication (the ray).  The alphabet of a known KS set can often be read directly from its geometric representation: in the cube diagrams used to display the Peres~33, Conway--Kochen~31, and related constructions~\cite{Peres1991,Peres1993}, each vertex carries explicit integer coordinates, and the set of distinct coordinate values appearing across all vertices is precisely the alphabet $\mathcal{A}$.

Two rays $v, w \in \mathcal{A}^3$ are orthogonal when their dot product vanishes: $\sum_k \bar v_k w_k = 0$.  For coordinates drawn from a finite alphabet, this requires an algebraic relation among elements of $\mathcal{A}$ that makes the sum cancel exactly---a \emph{cancellation identity}.  The cancellation identity is the mechanism by which an alphabet produces orthogonal rays, and hence the mechanism by which it can potentially generate a Kochen--Specker set.

For the integer alphabet $\{0, \pm 1, \pm 2\}$, the key identity is $1 + 1 = 2$: two coordinates of magnitude~1 can sum to cancel a single coordinate of magnitude~2, as in $1 \cdot 1 + 1 \cdot 1 + (-1) \cdot 2 = 0$.  For the Peres alphabet $\{0, \pm 1, \pm\sqrt{2}\}$, it is $(\sqrt{2})^2 = 2 = 1 + 1$.  For the Eisenstein alphabet $\{0, \pm 1, \pm\omega, \pm\bar\omega\}$ (where $\omega = e^{2\pi i/3}$), it is the three-term phase cancellation $1 + \omega + \omega^2 = 0$.  In every case, the identity has the same effect: it creates exact orthogonalities among triples of rays (triads), and when enough triads interlock, no consistent $\{0,1\}$-coloring exists---a KS set.

One might ask whether the trivial identity $1 + (-1) = 0$ already suffices---after all, it makes dot products vanish (e.g., $1 \cdot 1 + 1 \cdot (-1) + 0 \cdot 0 = 0$).  The answer is that the minimal alphabet $\{0, \pm 1\}$ generates only 13 projective rays in $\R^3$ (the 3 axis directions, 6 face diagonals, and 4 body diagonals of the cube), far too few for a KS set (the minimum known is 31).  The identity $1 + 1 = 2$ is what unlocks the $\pm 2$ coordinate, expanding the pool from 13 to 49 rays---nearly four times as many---and creating the dense triad network that KS-uncolorability requires.  In this sense, $1 + 1 = 2$ is the simplest \emph{useful} cancellation identity: the first one that generates enough orthogonal structure for the coloring constraints to become unsatisfiable.

The methodological approach of this paper is to \emph{reverse} this chain: we systematically construct coordinate alphabets from algebraic number fields, identify their cancellation identities, generate the induced ray sets, and test for KS-uncolorability.  This turns the search for new KS sets into a question of algebra: which rings admit cancellation identities of low enough norm to produce the dense orthogonality structures that KS-uncolorability requires?

\subsection{Cross-product completion}

\begin{definition}\label{def:completion}
The \emph{cross-product completion} of a ray set $S \subset \R^3$ is the closure $\bar{S}$ obtained by iterating: for every orthogonal pair $v, w \in S$, add the ray $[v \times w]$ to $S$, until no new rays are generated. This is well-defined since $v \perp w$ implies $v \times w \neq 0$ and $v \times w$ is orthogonal to both $v$ and $w$. Termination is not guaranteed in general but is observed empirically for all alphabets in our survey, typically stabilizing within 3--5 iterations (see Section~\ref{sec:limitations}).
\end{definition}

Cross-product completion is a natural choice of closure operation: it ensures that every orthogonal pair is extended to a complete basis, which is the minimal requirement for a ray set to carry KS-type coloring constraints.  Other closure operations are possible (e.g., closure under orthogonal complementation of arbitrary subsets, or Gram--Schmidt processes with field constraints), and the island classification may in principle depend on which closure is used.  We adopt cross-product completion throughout and note that it terminates empirically for all tested alphabets; a proof of termination for specific families of coordinate rings would strengthen the classification.  Since cross products are polynomial in the coordinates, completion trivially preserves any number field $K$. The nontrivial property is \emph{finiteness}: for a given alphabet $\mathcal{A}$, the completed ray pool $\bar{S}(\mathcal{A})$ stabilizes after finitely many iterations, and its size and triad structure depend sensitively on the algebraic identities available in $\mathcal{A}$.

\section{Computational Methods}\label{sec:methods}

\subsection{Ray canonicalization}\label{sec:canon}

A ray is an equivalence class of nonzero vectors under scalar multiplication. For \emph{real} alphabets ($\mathcal{A} \subset \R$), two vectors represent the same ray iff they are nonzero real multiples of each other; we choose the primitive representative with first nonzero entry positive. For \emph{complex} alphabets ($\mathcal{A} \subset \C$), two vectors represent the same ray iff they are nonzero complex multiples of each other (i.e., $v \sim \lambda v$ for any $\lambda \in \C^*$); we choose a canonical representative by dividing by the first nonzero coordinate (making it equal to 1), then applying a lexicographic tie-breaker on the remaining coordinates.

In both cases, representatives are \emph{unnormalized}: we do not divide by the vector's norm. This avoids introducing spurious irrationals through normalization and ensures that ray counts (e.g., 49, 109, 205) are well-defined projective counts over the coordinate ring. The canonicalization is deterministic and injective on rays: each projective equivalence class maps to exactly one representative. For real alphabets, the sign convention (first nonzero entry positive) resolves the $\pm v$ ambiguity; for complex alphabets, dividing by the first nonzero coordinate resolves the full $\C^*$ scaling ambiguity. Orthogonality is tested on the unnormalized representatives: $v \perp w$ iff $\sum_k \bar{v}_k w_k = 0$ (Hermitian, reducing to $\sum_k v_k w_k = 0$ in the real case). Since orthogonality is invariant under scaling, this test is independent of the choice of representative.

\subsection{KS coloring via SAT}

We encode the KS coloring problem as a Boolean satisfiability (SAT) instance. For a ray set with $n$ rays, orthogonal pairs $E$, and triads $T$:

\begin{itemize}
    \item Boolean variable $x_i$ indicates ray $i$ is green ($f(i) = 1$).
    \item For each triad $\{i, j, k\} \in T$: exactly-one constraint (4 clauses).
    \item For each pair $(i, j) \in E$: at-most-one constraint (1 clause: $\neg x_i \lor \neg x_j$).
\end{itemize}

We use the Glucose4 solver~\cite{AudemardSimon2018} via the PySAT library. UNSAT = KS-uncolorable. For a typical 33-ray, 16-triad configuration, the SAT instance has 33 variables and $4 \times 16 + |E| \approx 120$--$185$ clauses (depending on the number of orthogonal pairs $|E|$). No symmetry breaking is applied; the instances are small enough for direct solving (milliseconds). All SAT instances, ray coordinates, and basis lists are available in the accompanying code repository.

\subsection{Randomized greedy minimization}

Given a KS-uncolorable set $S$, we find minimal KS subsets by randomized greedy removal: shuffle the rays, attempt to remove each in random order, keeping the removal if the set remains KS-uncolorable. We run 200--1000 independent trials per configuration and report the smallest KS subset found.

\subsection{Realizability via numerical optimization}

Given an abstract orthogonality graph $G$ on $n$ vertices with edge set $E$, we test realizability in $\R^3$ by parametrizing each ray by spherical coordinates $(\theta_i, \phi_i)$ and minimizing
\[
\mathcal{L}(\theta, \phi) = \sum_{(i,j) \in E} (v_i \cdot v_j)^2
\]
using L-BFGS-B with 50 random restarts. We declare realizability if $\mathcal{L} < 10^{-8}$.

\section{Real Coordinate Alphabets}\label{sec:real}

\subsection{Systematic alphabet survey}

Table~\ref{tab:alphabets} summarizes the results of exhaustive ray generation and KS testing for coordinate alphabets based on various algebraic extensions of $\Q$.

\begin{table}[ht]
\centering
\caption{KS sets from real coordinate alphabets $\{0, \pm 1, \pm x\}$ in $\R^3$.}
\label{tab:alphabets}
\begin{tabular}{lcccccc}
\toprule
Alphabet & Rays & Pairs & Triads & KS? & Min \\
\midrule
$\{0, \pm 1\}$ & 13 & 24 & 4 & No & --- \\
$\{0, \pm 1, \pm\sqrt{2}\}$ & 49 & 120 & 16 & \textbf{Yes} & \textbf{33} \\
$\{0, \pm 1, \pm 2\}$ & 49 & 138 & 26 & \textbf{Yes} & \textbf{31} \\
$\{0, \pm 1, \pm\sqrt{3}\}$ & 49 & 114 & 10 & No & --- \\
$\{0, \pm 1, \pm\sqrt{5}\}$ & 49 & 114 & 10 & No & --- \\
$\{0, \pm 1, \pm\varphi\}$ & 49 & 138 & 10 & No & --- \\
$\{0, \pm 1, \pm\frac{1}{2}\}$ & 49 & 138 & 26 & \textbf{Yes} & \textbf{31} \\
\bottomrule
\end{tabular}
\end{table}

The Conway--Kochen 31-vector set is rediscovered by the integer alphabet in 19\% of randomized minimization trials (192/1000). No trial produced a set smaller than 31.

\subsection{The 2:1 ratio as the critical feature}\label{sec:ratio}

\begin{observation}
Every real alphabet that achieves a 31-vector KS set contains the ratio $2{:}1$ among its nonzero elements---either as $\pm 2$ alongside $\pm 1$, or equivalently as $\pm \frac{1}{2}$ alongside $\pm 1$. Adding other algebraic values (e.g., $\sqrt{2}$, $\varphi$, $\sqrt{3}$) to such an alphabet does not reduce the minimum below 31.
\end{observation}

Table~\ref{tab:extended} shows extended alphabets. In every case, the presence of the 2:1 ratio yields 31; its absence yields $\geq 33$ or colorable.

\begin{table}[ht]
\centering
\caption{Extended alphabets: across all tested configurations, the 2:1 ratio correlates perfectly with reaching 31.}
\label{tab:extended}
\begin{tabular}{lcccc}
\toprule
Alphabet & Rays & Triads & Min KS & Has 2:1? \\
\midrule
$\{0,\pm 1,\pm\sqrt{2},\pm 2\}$ & 109 & 38 & 31 & Yes \\
$\{0,\pm 1,\pm\varphi,\pm 2\}$ & 145 & 50 & 31 & Yes \\
$\{0,\pm 1,\pm 2,\pm 3\}$ & 145 & 50 & 31 & Yes \\
$\{0,\pm 1,\pm\sqrt{2},\pm\varphi\}$ & 145 & 28 & 33 & No \\
$\{0,\pm 1,\pm\varphi,\pm\varphi^2\}$ & 109 & 24 & colorable & No \\
\bottomrule
\end{tabular}
\end{table}

Among the two-element alphabets tested, $1 + 1 = 2$ produces the highest triad density: 26 triads from 49 rays. This density appears to be the mechanism by which the integer alphabet achieves the smallest minimum of 31. Heuristically, the identity $a + a = b$ with $a, b$ both in the alphabet generates a large number of three-term dot-product cancellations and their coordinate permutations. Identities involving irrationals (e.g., $(\sqrt{2})^2 = 2$) produce fewer cancellations because they require squaring to reach an integer, reducing the number of distinct triad-generating permutations. We note that ``unique'' here is relative to the alphabets in our survey and the equivalence class under scaling; a formal proof of optimality among all possible two-element cancellation identities remains open.

\subsection{Computational evidence that 31 is minimum for the integer island}\label{sec:exhaustive}

We provide strong computational evidence that no smaller KS set exists within the integer island's ray pool, though a complete proof over the full pool remains computationally infeasible.

Starting from the 49-ray pool of $\{0, \pm 1, \pm 2\}$, we ran 1{,}000 independent randomized greedy minimization trials, each producing a minimal KS subset. Six distinct 31-ray sets were found, collectively using exactly 37 of the 49 available rays. We then performed an exhaustive check: every $\binom{37}{30} = 10{,}295{,}472$ subset of size 30 drawn from these 37 rays was tested for KS-uncolorability via SAT. All were colorable.

\begin{proposition}\label{prop:31optimal}
No 30-ray subset of the 37-ray union of all discovered minimal 31-sets is KS-uncolorable. Furthermore, no single-ray removal, two-ray removal, or ray-swap operation applied to any of the six 31-sets produces a KS-uncolorable 30-ray configuration.
\end{proposition}

\begin{remark}
The exhaustive check above covers the 37-ray union of discovered 31-sets, not the entire 49-ray pool. We close this gap using an OCUS (Optimal Constrained Unsatisfiable Subset) approach: for each target size $n \leq 30$, we add a cardinality constraint $\sum_v b_v \leq n$ to the SAT encoding and use a CEGAR loop---when the solver finds a satisfying coloring, we extract a blocking clause and iterate until UNSAT (proving no $n$-ray KS subset exists) or a witness is found. This exhaustive procedure proves that no KS-uncolorable subset of size $\leq 30$ exists within the full 49-ray pool, for all $n$ from 21 to 30 (each size exhausted in under 1\,s; $n = 30$ required 272 CEGAR iterations). As independent confirmation, 1{,}000 MUS (Minimal Unsatisfiable Subset) extractions all returned subsets of size~31, yielding 162 distinct minimal KS sets.
\end{remark}

The exhaustive check in Proposition~\ref{prop:31optimal} required approximately 10.3 million SAT calls. As an independent verification, we also performed neighborhood searches around each of the six 31-sets: all 31 single-ray removals (per set), all 8{,}370 ray-swap operations, and all 465 two-ray removals per set were tested. Every resulting 30-ray configuration was colorable.

This result establishes that 31 is optimal \emph{within the full 49-ray integer pool}.  Combined with the rigidity theorem of Trandafir and Cabello~\cite{TrandafirCabello2025} (who proved CK-31 is unique up to unitary equivalence), and the absence of any sub-31 KS set across all six algebraic islands, this is consistent with the conjecture that 31 is optimal for all KS sets in $\R^3$---though it does not constitute a proof, as coordinates outside the integer alphabet are not excluded.

\subsection{Quadratic number fields}\label{sec:quadratic}

We systematically tested the alphabet $\{0, \pm 1, \pm\sqrt{d}\}$ for every non-square integer $d = 2, \ldots, 30$.

\begin{proposition}\label{prop:quadratic}
Among the quadratic fields $\Q(\sqrt{d})$ for $d \in \{2, 3, 5, 6, 7, \ldots, 30\}$ (excluding perfect squares), only $d = 2$ produces a KS-uncolorable set from the standard two-element alphabet $\{0, \pm 1, \pm\sqrt{d}\}$. All others have $\leq 10$ triads and are colorable.
\end{proposition}

A heuristic explanation: the identity $(\sqrt{d})^2 = d$ produces a three-term dot-product cancellation only when $d$ equals a sum of alphabet elements. For $d = 2$, we have $(\sqrt{2})^2 = 1 + 1$, giving the Peres cancellation. For $d \geq 3$, no comparable two-term sum from $\{1, \sqrt{d}\}$ equals $d$, so the cancellation structure is too sparse to generate sufficient triads. (A proof that no other cancellation pattern can arise for these alphabets would require a more detailed case analysis, which we leave to future work.)

\section{Complex Coordinate Alphabets}\label{sec:complex}

\subsection{Roots of unity and the $6 \mid n$ conjecture}\label{sec:roots}

We generated ray sets from the alphabet $\{0\} \cup \{\zeta^k : k = 0, \ldots, n-1\}$ where $\zeta = e^{2\pi i/n}$, using the Hermitian inner product $\braket{v}{w} = \sum_k \bar{v}_k w_k$ for orthogonality. Table~\ref{tab:roots} shows results for $n = 2, \ldots, 30$.

\begin{table}[ht]
\centering
\caption{KS sets from $n$th roots of unity. Only multiples of 6 (bold) are uncolorable. The ``Best found'' column reports the smallest KS subset found by heuristic minimization; these are upper bounds, not proven minima. For $6 \mid n$, subset containment guarantees $\leq 33$ (see text).}
\label{tab:roots}
\begin{tabular}{ccccccc}
\toprule
$n$ & $6 \mid n$? & Rays & Triads & KS? & Best found & $\leq 33$? \\
\midrule
2--5 & No & 13--43 & 1--7 & No & --- & --- \\
\textbf{6} & \textbf{Yes} & \textbf{57} & \textbf{22} & \textbf{Yes} & \textbf{33} & \textbf{Yes} \\
7--11 & No & 73--157 & 1--28 & No & --- & --- \\
\textbf{12} & \textbf{Yes} & \textbf{183} & \textbf{67} & \textbf{Yes} & \textbf{33} & \textbf{Yes} \\
13--17 & No & 211--343 & 1--76 & No & --- & --- \\
\textbf{18} & \textbf{Yes} & \textbf{381} & \textbf{136} & \textbf{Yes} & \textbf{76} & \textbf{Yes} \\
19--23 & No & 421--601 & 1--148 & No & --- & --- \\
\textbf{24} & \textbf{Yes} & \textbf{651} & \textbf{229} & \textbf{Yes} & --- & \textbf{Yes} \\
25--29 & No & 703--931 & 1--244 & No & --- & --- \\
\textbf{30} & \textbf{Yes} & \textbf{993} & \textbf{346} & \textbf{Yes} & --- & \textbf{Yes} \\
\bottomrule
\end{tabular}
\end{table}

\begin{conjecture}\label{conj:6n}
The ray set generated by $n$th roots of unity in $\C^3$ is KS-uncolorable if and only if $6 \mid n$.
\end{conjecture}

This has been verified computationally for all $2 \leq n \leq 30$ with no exceptions. A proof for general $n$ would likely follow from the representation theory of cyclic groups acting on $\C^3$; we sketch the necessary and sufficient conditions below but leave a complete proof to future work.

\begin{remark}[Subset containment]
For any $n$ divisible by 6, the $n$th roots of unity contain the 6th roots as a subset, so the $n=6$ ray pool embeds into the $n$-ray pool. Consequently, any $6\mid n$ pool contains a 33-vector KS subset (the Eisenstein-based set from $n=6$). The ``Best found'' column in Table~\ref{tab:roots} reports the output of heuristic greedy minimization, which for $n=18$ returned a 76-vector set---a failure of the minimizer, not a true lower bound. The true minimum for every $6\mid n$ pool is at most 33.
\end{remark}

The mechanism requires two ingredients working in concert:
\begin{itemize}
    \item \emph{Three-term phase cancellation} from $3 \mid n$: the identity $1 + \omega + \omega^2 = 0$ (where $\omega = e^{2\pi i/3}$) creates orthogonal triads.
    \item \emph{Real-axis orthogonalities} from $2 \mid n$: the real embedding ($\zeta^{n/2} = -1$) creates orthogonal pairs that interlock the triads.
\end{itemize}

Divisibility by 3 alone (e.g., $n = 9, 15, 21, 27$) creates many triads but their constraint network is satisfiable. Divisibility by 2 alone (e.g., $n = 4, 8, 10$) creates pairs but too few triads. Only when both are present ($6 \mid n$) does the interlocking force a KS contradiction.

\begin{remark}
This result connects to the finding of Cortez, Morales, and Reyes~\cite{CortezMoralesReyes2022} that $\Z[1/6]$ is the minimal ring extension of $\Z$ exhibiting KS contextuality in dimension 3. The prime factorization $6 = 2 \times 3$ appears in both results: the divisibility condition $6 \mid n$ in the cyclotomic setting mirrors the necessity of inverting both 2 and 3 in the ring-theoretic setting. We can make this connection quantitative using their $N(S) = \operatorname{lcm}\{\|v\|^2 : v \in S\}$ invariant, which determines the ring $\Z[1/N(S)]$ in which the projection matrices $P_v = vv^T/\|v\|^2$ live. Computing $N(S)$ for each island's minimal KS set: Eisenstein $N = 6$ ($= 2 \times 3$), Peres $N = 12$ ($= 2^2 \times 3$), $\Z[\sqrt{-2}]$ $N = 12$, Heegner-7 $N = 12$, CK-31 $N = 30$ ($= 2 \times 3 \times 5$). The Eisenstein island achieves exactly the Cortez--Morales--Reyes minimum $N = 6$, confirming it as the algebraic realization of their minimal ring. Every island requires primes 2 and~3; the integer island additionally needs~5 (from the $\|v\|^2 = 5$ rays in CK-31).
\end{remark}

\subsection{Eisenstein integers}

The Eisenstein integers $\Z[\omega]$ (where $\omega = e^{2\pi i/3}$) produce KS sets with smallest found subset of 33 vectors across all coefficient bounds tested (max coefficient 1--3, squared norm cutoff 3--7). Increasing the pool size beyond 57 rays does not reduce the smallest found below 33.

\subsection{Complex quadratic rings}\label{sec:complex-quadratic}

Extending the real quadratic survey to complex quadratic rings $\Z[\sqrt{-d}]$ with basic alphabets $\{0, \pm 1, \pm\sqrt{-d}\}$, we tested $d = 1, 2, \ldots, 15$. Most produce colorable ray sets. However, $d = 2$ yields a new island:

\begin{table}[ht]
\centering
\caption{Complex quadratic rings $\Z[\sqrt{-d}]$ with alphabet $\{0, \pm 1, \pm\sqrt{-d}\}$.}
\label{tab:complex-quadratic}
\begin{tabular}{lcccccc}
\toprule
Ring & $d$ & Rays & Pairs & Triads & KS? & Min \\
\midrule
$\Z[i]$ & 1 & 25 & 48 & 6 & No & --- \\
$\Z[\sqrt{-2}]$ & 2 & 49 & 120 & 16 & \textbf{Yes} & \textbf{33} \\
$\Z[\sqrt{-3}]$ & 3 & 49 & 114 & 10 & No & --- \\
$\Z[\sqrt{-5}]$ & 5 & 49 & 114 & 10 & No & --- \\
$\Z[\sqrt{-7}]$ & 7 & 49 & 114 & 10 & No & --- \\
\bottomrule
\end{tabular}
\end{table}

The ring $\Z[\sqrt{-2}]$ with alphabet $\{0, \pm 1, \pm\sqrt{-2}\}$ generates 49 rays, 120 orthogonal pairs, and 16 triads---an uncolorable configuration with smallest found KS subset of 33 vectors. The cancellation identity is $(\sqrt{-2})^2 = -2$, which in the Hermitian inner product yields $\overline{\sqrt{-2}} \cdot \sqrt{-2} = |\sqrt{-2}|^2 = 2$, producing the same numerical cancellation $1 + 1 = 2$ that drives the Peres island. Indeed, the ray count, pair count, and triad count of $\Z[\sqrt{-2}]$ exactly match those of $\Z[\sqrt{2}]$ (Table~\ref{tab:alphabets}), and the minimum KS subset size is identical (33). The two islands have isomorphic combinatorial structure despite living in different ambient fields ($\R$ vs.\ $\C$).

We also tested the ring of integers of $\Q(\sqrt{-7})$, namely $\Z[(1+\sqrt{-7})/2]$, with extended alphabet. This produced a larger pool (97 rays) that is KS-uncolorable with minimum found of 43 vectors---a sixth island with a larger minimum than the 33-vector cluster.

\section{Algebraic Island Classification}\label{sec:islands}

\subsection{Cross-product closure as a classification tool}\label{sec:closure}

An important subtlety is that KS behavior depends on the \emph{choice of generating alphabet}, not merely on the ambient number field. For instance, $\Q(\sqrt{5}) = \Q(\varphi)$ as fields, but the alphabet $\{0, \pm 1, \pm\sqrt{5}\}$ produces a colorable set while $\{0, \pm 1, \pm\varphi\}$ (where $\varphi = (1+\sqrt{5})/2$ is the ring-of-integers generator for $\Q(\sqrt{5})$) produces a KS-uncolorable set after completion. We therefore define:

\begin{definition}
An \emph{algebraic island} is a pair $(\mathcal{R}, \mathcal{A})$ consisting of a subring $\mathcal{R} \subseteq \C$ (typically the ring of integers of a number field) and a finite alphabet $\mathcal{A} \subset \mathcal{R}$ such that:
\begin{enumerate}[label=(\roman*)]
    \item the ray set $S(\mathcal{A})$ generates vectors in $\mathcal{R}^3$;
    \item the completion of $S(\mathcal{A})$ (see below) terminates in a \emph{finite} ray pool $\bar{S}$ with coordinates in $\mathcal{R}$;
    \item the completed ray pool $\bar{S}$ is KS-uncolorable.
\end{enumerate}
Two islands $(\mathcal{R}_1, \mathcal{A}_1)$ and $(\mathcal{R}_2, \mathcal{A}_2)$ are \emph{equivalent} if their completed ray pools have isomorphic orthogonality hypergraphs.
\end{definition}

The \emph{completion} in condition (ii) depends on whether the island is real or complex:

\begin{itemize}
    \item \textbf{Real islands} ($\mathcal{R} \subset \R$): cross-product completion (Definition~\ref{def:completion}). For each orthogonal pair $v \perp w$, adjoin $[v \times w]$.
    \item \textbf{Complex islands} ($\mathcal{R} \subset \C$, $\mathcal{R} \not\subset \R$): for each Hermitian-orthogonal pair $v \perp w$ in $\C^3$, adjoin the unique ray $[u]$ spanning $(\mathrm{span}\{v,w\})^\perp$. Concretely, $u$ is the vector of $2 \times 2$ cofactors of the matrix $\begin{pmatrix} \bar{v} \\ \bar{w} \end{pmatrix}$:
    \[
    u_k = (-1)^{k+1} \det \begin{pmatrix} \bar{v}_{k+1} & \bar{v}_{k+2} \\ \bar{w}_{k+1} & \bar{w}_{k+2} \end{pmatrix}, \quad k = 1,2,3 \pmod{3}.
    \]
    This is the Hermitian analogue of the cross product. When $\mathcal{R}$ is closed under conjugation, $u \in \mathcal{R}^3$, so completion preserves the coordinate ring.
\end{itemize}

\noindent In $\R^3$, conjugation is trivial and the cofactor formula reduces to the standard cross product, so the two definitions agree.

Note that condition (ii) is automatically satisfied when $\mathcal{R}$ is closed under the coordinate operations of the completion (cross products preserve $\mathcal{R}^3$ for real rings; orthogonal complements preserve $\mathcal{R}^3$ for complex rings closed under conjugation). The \emph{finiteness} of the completed pool and the \emph{combinatorial structure} it generates are the nontrivial properties. The same ring $\mathcal{R}$ may support zero, one, or multiple islands depending on the alphabet.

\begin{table}[ht]
\centering
\caption{Cross-product closure analysis for \emph{real} islands ($\mathcal{R} \subset \R$). ``Compl.''\ = completed ray count; ``New mag.''\ = distinct new coordinate magnitudes introduced by completion (e.g., for integers, cross products introduce 3, 4, 5 beyond $\{0,1,2\}$); ``Min'' = smallest KS subset found (200--1000 trials).}
\label{tab:closure}
\small
\begin{tabular}{lcccccc}
\toprule
Alphabet $(\mathcal{R}, \mathcal{A})$ & Raw & Compl. & New mag. & $2{:}1$? & Uncol.? & Min \\
\midrule
$(\Z, \{0,\pm 1,\pm 2\})$ & 49 & 109 & $+3$ & Yes & Yes & \textbf{31} \\
$(\Z[\sqrt{2}], \{0,\pm 1,\pm\!\sqrt{2}\})$ & 49 & 145 & $+12$ & No & Yes & \textbf{33} \\
$(\Z[\sqrt{3}], \{0,\pm 1,\pm\!\sqrt{3}\})$ & 49 & 145 & $+12$ & No & No & --- \\
$(\Z[\sqrt{5}], \{0,\pm 1,\pm\!\sqrt{5}\})$ & 49 & 145 & $+12$ & No & No & --- \\
$(\Z[\varphi], \{0,\pm 1,\pm\varphi\})$ & 49 & 205 & $+21$ & No & Yes & \textbf{52} \\
\bottomrule
\end{tabular}

\medskip
\noindent For the \emph{complex} islands, completion uses the Hermitian orthogonal-complement operation (Section~\ref{sec:closure}). The Eisenstein alphabet $\{0, \pm 1, \pm\omega, \pm\bar\omega\}$ with $\omega = e^{2\pi i/3}$ generates 57 rays (22 triads) at max coefficient 1. Complex completion does not expand this pool (all Hermitian-orthogonal complements already lie within it), and the completed set is KS-uncolorable with smallest found subset of 33 vectors. The complex quadratic alphabet $\{0, \pm 1, \pm\sqrt{-2}\}$ generates 49 rays (16 triads) with identical combinatorial structure to the Peres island; its smallest found KS subset is also 33 vectors. The Heegner-7 alphabet $\{0, \pm 1, \pm\alpha, \pm\bar\alpha\}$ with $\alpha = (1+\sqrt{-7})/2$ generates 145 rays (42 triads); its smallest found KS subset is 43 vectors (Section~\ref{sec:heegner}).
\end{table}

\subsection{Mixed alphabet evidence}\label{sec:mixed}

Cross-product completion of mixed alphabets---combining rays from two algebraically independent coordinate alphabets, then closing under cross products---provides further evidence for the primacy of the cancellation identity. The cross products generate rays in \emph{neither} original alphabet, exploring configurations invisible to any single-alphabet search.

\begin{table}[ht]
\centering
\caption{Mixed alphabet search with cross-product completion.}
\label{tab:mixed}
\begin{tabular}{lcccc}
\toprule
Combination & Total rays & Triads & KS? & Min \\
\midrule
$\{0,\pm 1,\pm 2\} + \{0,\pm 1,\pm\sqrt{2}\}$ & 205 & 158 & Yes & 31 \\
$\{0,\pm 1,\pm 2\} + \{0,\pm 1,\pm\sqrt{3}\}$ & 217 & 164 & Yes & 31 \\
$\{0,\pm 1,\pm 2\} + \{0,\pm 1,\pm\sqrt{5}\}$ & 217 & 164 & Yes & 31 \\
$\{0,\pm 1,\pm 2\} + \{0,\pm 1,\pm\varphi\}$ & 241 & 188 & Yes & 34$^\dagger$ \\
$\{0,\pm 1,\pm\sqrt{2}\} + \{0,\pm 1,\pm\sqrt{3}\}$ & 229 & 166 & Yes & 33 \\
$\{0,\pm 1,\pm\sqrt{2}\} + \{0,\pm 1,\pm\varphi\}$ & 253 & 190 & Yes & 33 \\
\bottomrule
\end{tabular}
\end{table}

Cross-product completion adds 150+ rays beyond the original alphabets, but the minimum KS size does not improve. The 31-vector barrier persists across all algebraically enriched constructions containing the 2:1 ratio. Conversely, combinations lacking the 2:1 ratio achieve at best 33 despite their larger pools, confirming that the 2:1 ratio is the rate-limiting algebraic feature, not the pool size.

\smallskip\noindent$^\dagger$The pool $\{0,\pm 1,\pm 2\} + \{0,\pm 1,\pm\varphi\}$ contains all 49 integer-alphabet rays as a subset, so it includes CK-31 as a sub-configuration. The reported minimum of 34 reflects the fact that our randomized greedy minimization (200 trials) did not locate the 31-ray subset within the larger 241-ray pool.

\begin{remark}[Relation to the SI-C closure of Trandafir and Cabello]
Trandafir and Cabello~\cite{TrandafirCabello2025} construct their KS sets by starting from the 13-element minimal state-independent contextuality (SI-C) set of Yu and Oh, completing to orthonormal bases (25 rays), and adding all vectors orthogonal to $\geq 2$ existing vectors (97 rays total).  We computed the overlap between this 97-ray SI-C closure and our 49-ray integer pool $\{0,\pm1,\pm2\}$: they share 37 rays, with 60 rays in the SI-C closure only (coordinates up to $\pm 5$, norms up to~35) and 12 rays in the integer pool only (those with two $\pm 2$ coordinates, norm~9).  A second round of orthogonal closure grows the SI-C set to 1{,}741~rays that \emph{contain} the entire integer pool as a subset.  The two constructions---their top-down closure from SI-C and our bottom-up alphabet enumeration---thus converge: CK-31 lives in the intersection of both frameworks.
\end{remark}

\subsection{The six known islands}\label{sec:sixislands}

\begin{observation}\label{obs:islands}
Among all alphabets and fields tested in our survey---quadratic fields $\Q(\sqrt{d})$ for $d = 2, \ldots, 30$, cyclotomic fields for $n \leq 30$, and selected extensions---six algebraic islands support KS sets:
\begin{enumerate}[label=(\arabic*)]
    \item \textbf{Integer island} $(\Z, \{0,\pm 1,\pm 2\})$: cancellation $1 + 1 = 2$; smallest found: 31 vectors ($\CK$). Supported by exhaustive enumeration over the 37-ray union of discovered 31-sets (Proposition~\ref{prop:31optimal}).
    \item \textbf{Peres island} $(\Z[\sqrt{2}], \{0,\pm 1,\pm\sqrt{2}\})$: cancellation $(\sqrt{2})^2 = 1 + 1$; smallest found: 33 vectors.
    \item \textbf{Eisenstein island} $(\Z[\omega], \{0,\pm 1,\pm\omega,\pm\bar\omega\})$: cancellation $1 + \omega + \omega^2 = 0$; smallest found: 33 vectors. Identified with Cabello's ``simplest KS set''~\cite{Cabello2025simplest} (Section~\ref{sec:cabello}).
    \item \textbf{Complex quadratic island} $(\Z[\sqrt{-2}], \{0,\pm 1,\pm\sqrt{-2}\})$: cancellation $|\sqrt{-2}|^2 = 2$; smallest found: 33 vectors. A new algebraic realization of the norm-2 cancellation; graph-isomorphic to the Peres island (Proposition~\ref{prop:isomorphism}).
    \item \textbf{Heegner-7 island} $(\Z[(1+\sqrt{-7})/2], \{0,\pm 1,\pm\alpha,\pm\bar\alpha\})$ where $\alpha = (1+\sqrt{-7})/2$: cancellation $\alpha\bar\alpha = 2$; smallest found: 43 vectors. Newly discovered (Section~\ref{sec:heegner}).
    \item \textbf{Golden island} $(\Z[\varphi], \{0,\pm 1,\pm\varphi\})$: cancellation $\varphi^2 = \varphi + 1$; smallest found: 52 vectors.
\end{enumerate}
No other alphabet from our survey (real quadratic fields $\Q(\sqrt{d})$ for $d = 2,\ldots,30$; all nine imaginary quadratic fields with class number 1; roots of unity with $6\nmid n$; cubic extensions; icosahedral symmetry directions) produced a KS-uncolorable set. The hierarchy $31 < 33 = 33 = 33 < 43 < 52$ was consistent across 200--1000 randomized minimization trials per configuration; the size distributions are reported in the supplementary data.
\end{observation}

This hierarchy correlates with the algebraic complexity of the cancellation identity:

\begin{table}[ht]
\centering
\caption{The efficiency hierarchy: simpler cancellations yield smaller KS sets. ``Smallest found'' reports the best result from randomized greedy minimization (200--1000 trials).}
\label{tab:hierarchy}
\begin{tabular}{clcccc}
\toprule
Rank & Identity & Distinct magnitudes & Triads (raw) & Smallest found \\
\midrule
1 & $1 + 1 = 2$ & 2 & 26 & \textbf{31} \\
2 & $(\sqrt{2})^2 = 1 + 1$ & 2 (after squaring) & 16 & \textbf{33} \\
2$'$ & $1 + \omega + \omega^2 = 0$ & 1 (complex) & 22 & \textbf{33} \\
2$''$ & $|\sqrt{-2}|^2 = 2$ & 2 (complex) & 16 & \textbf{33} \\
3 & $\alpha\bar\alpha = 2$, $\alpha^2 - \alpha + 2 = 0$ & 2 (complex) & 42 & \textbf{43} \\
4 & $\varphi^2 = \varphi + 1$ & 3 & 10 & \textbf{52} \\
\bottomrule
\end{tabular}
\end{table}

\subsection{Discovery of the golden ratio island}\label{sec:golden}

The golden ratio field $\Q(\varphi)$ where $\varphi = (1 + \sqrt{5})/2$ presents a remarkable case. The raw alphabet $\{0, \pm 1, \pm\varphi\}$ generates 49 rays with only 10 triads---colorable (the SAT instance is satisfiable). However, cross-product completion generates 156 additional rays (all with coordinates in $\Q(\varphi)$), expanding the pool to 205 rays with 166 triads. This completed set \emph{is} KS-uncolorable, with the smallest KS subset found containing 52 vectors.

Note that $\varphi = (1+\sqrt{5})/2$ is the ring-of-integers generator for $\Q(\sqrt{5})$ (since $5 \equiv 1 \pmod{4}$). The alternative alphabet $\{0, \pm 1, \pm\sqrt{5}\}$ (using the field generator rather than the ring-of-integers generator) produces a \emph{different} and colorable set. This demonstrates that the choice of alphabet within a fixed number field is consequential (see Section~\ref{sec:closure}).

\begin{observation}\label{obs:golden}
The golden ratio island is invisible to standard alphabet searches because the raw alphabet $\{0, \pm 1, \pm\varphi\}$ is colorable. It becomes KS-uncolorable only through the emergent orthogonalities created by cross-product completion, which generates additional rays with coordinates in $\Z[\varphi]$ (the ring of integers of $\Q(\sqrt{5})$). This suggests there may be other algebraic substrates whose KS structure is hidden from raw-alphabet searches and revealed only by completion---though our survey of quadratic and selected higher extensions found no further examples.
\end{observation}

\subsection{Discovery of the Heegner-7 island}\label{sec:heegner}

A systematic search over the imaginary quadratic fields $\Q(\sqrt{-d})$ with class number 1---the \emph{Heegner numbers} $d = 1, 2, 3, 7, 11, 19, 43, 67, 163$---reveals a clean classification. For $d \equiv 3 \pmod{4}$, the ring of integers is $\Z[(1+\sqrt{-d})/2]$, not $\Z[\sqrt{-d}]$, and the generator $\alpha = (1+\sqrt{-d})/2$ satisfies $\alpha\bar\alpha = (1+d)/4$. The results:

\begin{table}[ht]
\centering
\caption{Imaginary quadratic fields with class number 1 (Heegner numbers). ``Gen.\ norm'' $= |\alpha|^2$ where $\alpha$ is the ring-of-integers generator. Only $d = 1, 2, 3, 7$ produce KS sets.}
\label{tab:heegner}
\begin{tabular}{cclcccl}
\toprule
$d$ & Gen.\ norm & Ring & Rays & Triads & Uncol.? & Min KS \\
\midrule
1 & 2 & $\Z[i]+(1+i)$ & 127 & 51 & Yes & 33 \\
2 & 2 & $\Z[\sqrt{-2}]$ & 49 & 16 & Yes & 33 \\
3 & 1 & $\Z[\omega]$ & 57 & 22 & Yes & 33 \\
\textbf{7} & \textbf{2} & $\Z[(1+\sqrt{-7})/2]$ & \textbf{145} & \textbf{42} & \textbf{Yes} & \textbf{43} \\
11 & 3 & $\Z[(1+\sqrt{-11})/2]$ & 145 & 30 & No & --- \\
19 & 5 & $\Z[(1+\sqrt{-19})/2]$ & 145 & 30 & No & --- \\
43 & 11 & $\Z[(1+\sqrt{-43})/2]$ & 145 & 30 & No & --- \\
67 & 17 & $\Z[(1+\sqrt{-67})/2]$ & 145 & 30 & No & --- \\
163 & 41 & $\Z[(1+\sqrt{-163})/2]$ & 145 & 30 & No & --- \\
\bottomrule
\end{tabular}
\end{table}

The pattern is sharp: KS-uncolorability requires the alphabet to support a cancellation identity whose terms have squared modulus at most~2. When $|\alpha|^2 = 2$, the identity $\alpha\bar\alpha = 2$ provides the same norm-2 cancellation as the integer $1+1=2$ and the Peres $(\sqrt{2})^2 = 2$, but through the complex Hermitian inner product. The Eisenstein case ($d = 3$) has $|\omega|^2 = 1$ but compensates with the three-term phase identity $1 + \omega + \omega^2 = 0$. At $|\alpha|^2 \geq 3$ ($d \geq 11$), the squared moduli are too large for dot-product cancellations in $\C^3$.

The Heegner-7 island is structurally distinct from all other islands. To our knowledge, neither the 43-vector Heegner-7 configuration nor the 52-vector golden ratio configuration has been previously reported in the KS literature; existing systematic searches~\cite{Pavicic2004, LiBrightGanesh2024} have focused on the range up to 31 vectors (see Appendix~\ref{app:rays} for explicit descriptions). The Heegner-7 minimal KS set has 43 vectors arranged in 23 bases, with degree distribution $\{4{:}16, 5{:}20, 6{:}4, 8{:}2, 10{:}1\}$ (five distinct degree types and 12 orbit types). This is richer than both the Eisenstein 33-set (three degree types, 5 orbits) and the Peres/complex-quadratic 33-sets (two degree types, 5--6 orbits), placing Heegner-7 between the 33-cluster and the Golden island in structural complexity.

\begin{observation}\label{obs:gaussian}
The Gaussian integers $\Z[i]$ with the enriched alphabet $\{0, \pm 1, \pm i, \pm(1+i)\}$ are also KS-uncolorable (min = 33), but the minimal set has degree distribution $\{4{:}30, 8{:}3\}$---identical to the Peres and $\Z[\sqrt{-2}]$ islands. This confirms that the classification is governed by the cancellation identity (here $|1+i|^2 = 2$), not the ambient ring: all three are realizations of the ``norm-2 cancellation'' and produce minimal KS sets sharing all tested graph invariants (Proposition~\ref{prop:isomorphism}).
\end{observation}

\begin{proposition}\label{prop:isomorphism}
\textbf{(Strong evidence for graph isomorphism of norm-2 islands.)} The Peres $(\Z[\sqrt{2}])$ and complex quadratic $(\Z[\sqrt{-2}])$ orthogonality graphs share all tested graph invariants, both at the full pool level (49 vertices, 120 edges) and at the minimized 33-set level (33 vertices, 72 edges): identical degree sequences, eigenvalue spectra $($all 49 eigenvalues match to $10^{-6})$, triangle counts $(16)$, and traces of $A^k$ through $k = 4$. The Weisfeiler--Leman 1-dimensional color refinement algorithm stabilizes at 6 color classes $($full pool$)$ and 3 color classes $($minimized$)$ without distinguishing the graphs.  These invariants provide strong evidence for isomorphism but do not constitute a certified proof; a definitive confirmation would require a graph isomorphism solver such as nauty/traces providing an explicit vertex bijection.
\end{proposition}

This reduces the six known algebraic islands to five distinct orthogonality graph types. The Peres, complex quadratic, and Gaussian islands share a single graph type, unified by the norm-2 cancellation identity $|\alpha|^2 = 2$ operating in different rings ($\Z[\sqrt{2}]$, $\Z[\sqrt{-2}]$,~$\Z[i]$).

\subsection{Dead ends}

\paragraph{Icosahedral symmetry.} The icosahedron has 31 symmetry-axis directions---6 vertex, 15 edge, 10 face---a tantalizing coincidence with $\CK$. However, these directions produce only 5 triads (colorable). After cross-product completion: 91 rays, 65 triads, still colorable. A heuristic explanation: the icosahedral group has rotation orders 2, 3, 5 (angles 180$^\circ$, 120$^\circ$, 72$^\circ$) but \emph{no 90$^\circ$ rotations}, whereas the orthogonal group of $\CK$ is octahedral, the largest finite subgroup of SO(3) containing 90$^\circ$ rotations.

\paragraph{Higher quadratic fields.} All $\Q(\sqrt{d})$ for $d = 3, \ldots, 30$ (excluding $d = 4, 9, 16, 25$) produce $\leq 10$ triads and are colorable, even after cross-product completion (for $d \neq 5$; $d = 5$ generates the golden ratio island as $\Q(\sqrt{5}) = \Q(\varphi)$).

\section{Discussion}\label{sec:discussion}

\subsection{Connections to ring-theoretic results}

The $6 \mid n$ conjecture (Conjecture~\ref{conj:6n}) and the Cortez--Morales--Reyes result that $\Z[1/6]$ is the minimal KS ring~\cite{CortezMoralesReyes2022} both identify the primes 2 and 3 as the arithmetic backbone of three-dimensional contextuality. In the cyclotomic setting, divisibility by 2 provides real-axis orthogonalities and divisibility by 3 provides three-term phase cancellations; both are necessary and jointly sufficient for KS-uncolorability. In the ring-theoretic setting, inverting 2 and 3 is necessary and sufficient for the existence of a KS contextual set of vectors with entries in $\Z[1/N]$. The precise relationship between these two formulations deserves further investigation.

\subsection{Why six islands?}

Among all tested coordinate alphabets, we observe six algebraic islands, arising from two fundamentally different cancellation mechanisms:

\begin{enumerate}
    \item[\textbf{(A)}] \textbf{Norm-2 cancellation}: an algebraic identity producing the value~2 from products of ring elements, enabling three-term dot-product cancellations of the form $a_1 b_1 + a_2 b_2 + a_3 b_3 = 0$ with bounded coordinates. Five islands use this mechanism:
    \begin{itemize}
        \item \textbf{Integer} ($1 + 1 = 2$): additive doubling of units.
        \item \textbf{Peres} ($(\sqrt{2})^2 = 2$), \textbf{Complex quadratic} ($|\sqrt{-2}|^2 = 2$), and \textbf{Gaussian} ($|1+i|^2 = 2$): squaring (or modulus-squaring) produces 2. These three share the same minimal combinatorial structure (degree distribution $\{4{:}30, 8{:}3\}$, min 33) despite living in different rings.
        \item \textbf{Heegner-7} ($\alpha\bar\alpha = 2$ where $\alpha = (1+\sqrt{-7})/2$): norm-2 cancellation through a quadratic integer satisfying $\alpha^2 - \alpha + 2 = 0$. Despite sharing the norm-2 identity with the Peres family, the richer alphabet (7 elements vs.\ 5) produces a more complex orthogonality graph (42 triads from 145 rays), yielding a larger minimum KS set of 43 vectors.
    \end{itemize}
    \item[\textbf{(B)}] \textbf{Phase cancellation}: a root-of-unity identity $1 + \omega + \omega^2 = 0$ (where $\omega = e^{2\pi i/3}$), in which three complex unit-norm elements sum to zero. This mechanism produces a qualitatively different three-term cancellation: the vanishing sum involves three equal-magnitude terms rather than two small terms balancing one larger term. The \textbf{Eisenstein island} ($\Z[\omega]$, min 33) is the sole representative.
\end{enumerate}

The \textbf{golden ratio island} ($\varphi^2 = \varphi + 1$, min 52) requires special comment. The generator $\varphi$ has $|\varphi|^2 \approx 2.618 > 2$, so it does not participate directly in a norm-2 cancellation identity.  However, cross-product completion generates the reciprocal $1/\varphi$ (with $|1/\varphi|^2 < 1$), and the completed coordinate set contains the ratio $\varphi : 1/\varphi = \varphi^2 = \varphi + 1$, reintroducing norm-2-type cancellations indirectly. The golden island demonstrates that the relevant criterion is the cancellation structure of the \emph{completed} coordinate set, not the raw alphabet: it is possible for cross-product completion to introduce cancellation mechanisms absent in the original alphabet.  This is also why the two-element alphabet classification (Observation~\ref{obs:two-element}) applies only to raw alphabets before completion.

Each mechanism generates a distinct orthogonality structure in its completed ray set, and the efficiency of the mechanism (measured by the minimum KS set size) correlates inversely with its algebraic complexity. The integer identity involves two distinct magnitudes and produces the densest triad network (26 triads from 49 rays, ratio 0.53). The golden ratio identity involves three distinct magnitudes and produces the sparsest (10 triads from 49 rays, ratio 0.20), requiring completion to 205 rays before becoming KS-uncolorable.

The apparent isomorphism between the Peres ($\Z[\sqrt{2}]$), complex quadratic ($\Z[\sqrt{-2}]$), and Gaussian ($\Z[i]$ with $1+i$) islands is noteworthy: three different rings produce minimal KS sets with identical graph invariants because the Hermitian inner product reduces $\overline{\sqrt{-2}} \cdot \sqrt{-2}$, $(\sqrt{2})^2$, and $\overline{(1+i)} \cdot (1+i)$ all to the same value 2. This confirms that the ``true'' classification is by cancellation identity, not by ambient field.  This phenomenon was first observed by Gould and Aravind~\cite{GouldAravind2010}, who showed that the Peres and Penrose 33-vector sets share the same orthogonality diagram and belong to a continuous three-parameter family with constraint $|c|^2 = 2$---precisely the norm-2 condition.  That result provides the analytic explanation for why norm-2 cancellation controls graph isomorphism across our islands.

The Heegner-7 island is especially significant because it demonstrates that the norm-2 identity can produce \emph{different} minimal KS set sizes depending on the generator's minimal polynomial. The generator $\alpha = (1+\sqrt{-7})/2$ satisfies $\alpha^2 = \alpha - 2$, introducing a linear term absent in the Peres case ($(\sqrt{2})^2 = 2$). This additional algebraic structure enriches the ray pool (145 vs.\ 49 rays) while simultaneously making the orthogonality graph harder to ``compress'' (min 43 vs.\ 33).

\subsection{Cancellation classification for two-element alphabets}\label{sec:completeness}

Table~\ref{tab:cancellation} summarizes the cancellation identity, alphabet, and a concrete dot-product example for each island.

\begin{table*}[t]
\centering
\caption{The six algebraic islands: alphabets, cancellation identities, and concrete orthogonality examples. Each cancellation identity is the algebraic equation that enables the dot product to vanish; the example shows two specific rays whose orthogonality depends on that identity.}
\label{tab:cancellation}
\footnotesize
\begin{tabular}{@{}llllcr@{}}
\toprule
Island & Alphabet $\mathcal{A}$ & Identity & Dot-product example & Type & Min \\
\midrule
Integer & $\{0,\pm 1,\pm 2\}$ & $1 + 1 = 2$ & $1{\cdot}1 + 1{\cdot}1 - 1{\cdot}2 = 0$ & Norm-2 & 31 \\
Peres & $\{0,\pm 1,\pm\!\sqrt{2}\}$ & $(\sqrt{2})^2 = 2$ & $1{\cdot}1 + 1{\cdot}1 - \sqrt{2}{\cdot}\sqrt{2} = 0$ & Norm-2 & 33 \\
$\Z[\sqrt{-2}]$ & $\{0,\pm 1,\pm\!\sqrt{-2}\}$ & $|\sqrt{-2}|^2 = 2$ & $\bar 1{\cdot}1 + \bar 1{\cdot}1 + \overline{\sqrt{-2}}{\cdot}\sqrt{-2} = 0$ & Norm-2 & 33 \\
Gaussian & $\{0,\pm 1,\pm i,\pm(1{+}i)\}$ & $|1{+}i|^2 = 2$ & $\bar 1{\cdot}1 + \bar 1{\cdot}1 + \overline{(1{+}i)}{\cdot}(1{+}i) = 0$ & Norm-2 & 33 \\
Eisenstein & $\{0,\pm 1,\pm\omega,\pm\bar\omega\}$ & $1 + \omega + \omega^2 = 0$ & $1{\cdot}1 + 1{\cdot}\omega + 1{\cdot}\omega^2 = 0$ & Phase & 33 \\
Heegner-7 & $\{0,\pm 1,\pm\alpha,\pm\bar\alpha\}$ & $\alpha\bar\alpha = 2$ & $\bar 1{\cdot}1 + \bar 1{\cdot}1 - \bar\alpha{\cdot}\alpha = 0$ & Norm-2 & 43 \\
Golden & $\{0,\pm 1,\pm\varphi\} \to$ compl. & $\varphi^2 = \varphi + 1$ & $1{\cdot}1 + 1{\cdot}\varphi - 1{\cdot}\varphi^2 = 0$ & Indirect & 52 \\
\bottomrule
\end{tabular}
\end{table*}

We can classify all \emph{primitive three-term cancellation identities} available to two-element alphabets.  Consider an arbitrary alphabet $\{0, \pm 1, \pm x\}$ with $x \notin \{0, \pm 1\}$.  The Hermitian product set $\bar{\mathcal{A}} \cdot \mathcal{A}$ consists of $\{0, \pm 1, \pm x, \pm \bar{x}, \pm |x|^2\}$ (with $\bar{x} = x$ in the real case).  A non-trivial orthogonality $\sum_{k=1}^3 \bar{u}_k v_k = 0$ requires three elements of this product set to sum to zero, with at least one involving $x$ (otherwise the identity is trivial).  Table~\ref{tab:two-element} enumerates all such zero-sums.  Note that this classification applies to the \emph{raw} alphabet before completion; cross-product completion can introduce coordinates outside $\mathcal{A}$, potentially enabling additional cancellation identities in the completed pool (as occurs in the golden ratio case).

\begin{table}[h]
\centering
\caption{All primitive three-term zero-sums from the product set $\bar{\mathcal{A}} \cdot \mathcal{A}$ of a two-element alphabet $\{0, \pm 1, \pm x\}$. Each row is an algebraic constraint on~$x$; every solution corresponds to a known island. This classification applies to the raw alphabet before cross-product completion.}
\label{tab:two-element}
\small
\setlength{\tabcolsep}{4pt}
\begin{tabular}{@{}llll@{}}
\toprule
Zero-sum pattern & Identity on $x$ & Solution & Island \\
\midrule
$1 + 1 - x = 0$ & $x = 2$ & $x = 2$ & Integer \\
$1 + 1 - |x|^2 = 0$ & $|x|^2 = 2$ & $\sqrt{2},\, \sqrt{-2},\, \tfrac{1{+}\sqrt{-7}}{2},\, \ldots$ & Peres/$\Z[\!\sqrt{-2}]$/Heeg.-7 \\
$1 + x - |x|^2 = 0$ & $|x|^2 = 1 + x$ & $x = \varphi$ (real) & Golden \\
$1 + x + |x|^2\!\cdot\! x = 0$ & $x^2 + x + 1 = 0$ & $x = \omega = e^{2\pi i/3}$ & Eisenstein \\
$|x|^2 + |x|^2 - 1 = 0$ & $|x|^2 = \tfrac{1}{2}$ & $x = 1/\sqrt{2}$ & $\equiv$ Peres (rescaled) \\
$x + x - |x|^2 = 0$ & $|x|^2 = 2x$ & $x = 2$ (real) & $\equiv$ Integer \\
\bottomrule
\end{tabular}
\end{table}

\begin{observation}\label{obs:two-element}
For any two-element alphabet $\{0, \pm 1, \pm x\}$, every primitive three-term cancellation identity in the product set $\bar{\mathcal{A}} \cdot \mathcal{A}$ reduces to one of the six known islands (up to scaling equivalence).  This classification exhausts the cancellation identities available to two-element alphabets \emph{before completion}; it does not rule out the possibility that cross-product completion could introduce qualitatively new cancellation mechanisms in the completed pool, though no such case was observed in our survey.
\end{observation}

The completeness does \emph{not} extend to larger alphabets.  A three-element alphabet $\{0, \pm 1, \pm x, \pm y\}$ introduces cross-products $\pm xy$, $\pm \bar{x}y$, etc., enabling cancellation identities such as $1 + xy - |x|^2 = 0$ or $x + y + 1 = 0$ that involve two independent generators.  While our survey has tested many such alphabets (Table~\ref{tab:extended} and Section~\ref{sec:mixed}), the space of multi-element alphabets---particularly those arising from higher-degree number fields with multiple independent generators---is not exhaustively covered.  This is the principal remaining gap in the classification.

\begin{remark}[Trigonometric consistency check]\label{rem:trig}
As an independent test of the island classification, we swept the one-parameter family $\mathcal{A}(\theta) = \{0, \pm 1, \pm\cos\theta, \pm\sin\theta\}$ across 164 angles spanning $(0, \pi/2]$. The Pythagorean identity $\cos^2\theta + \sin^2\theta = 1$ serves as the cancellation mechanism. Exactly 5 distinct alphabets (at isolated angles: $\arctan(1/2)$, $\pi/6$, $\arctan(1/\varphi)$, $\arctan(1/\sqrt{2})$, $\pi/4$) are KS-uncolorable; all reduce to known islands via the norm-2 mechanism or golden ratio completion. No new islands emerge. Moving $1^\circ$ from any uncolorable angle destroys the property, consistent with the algebraic nature of the norm constraint.
\end{remark}

\subsection{Connection to perfect quantum strategies}\label{sec:cabello}

Our algebraic island framework gains new significance from three recent results of Cabello and collaborators that establish KS sets as the fundamental engine of \emph{perfect quantum strategies}.

Cabello~\cite{Cabello2024bipartite} proved that every bipartite perfect quantum strategy---a strategy achieving the maximum algebraic violation of a Bell inequality---defines a KS set. More precisely, the measurements achieving a perfect strategy on a Bell scenario must contain a KS-uncolorable subset. This means KS sets are not merely foundational curiosities but are \emph{operationally necessary} for the strongest form of quantum advantage in nonlocal games.

The connection deepens: Liu et al.~\cite{Cabello2024equivalences} established a chain of equivalences showing that face nonsignaling correlations, full Bell nonlocality, all-versus-nothing proofs, and pseudotelepathy are all manifestations of the same underlying structure---which, by the first result, is always a KS set. The algebraic islands we classify are therefore not just the algebraic substrates of KS sets; they are the algebraic substrates of \emph{all} perfect quantum strategies in dimension 3.

Most directly relevant to our framework, Cabello~\cite{Cabello2025simplest} constructed what he calls the ``simplest KS set'': a 33-vector, 14-basis configuration in $\C^3$ built from the Eisenstein integers $\Z[\omega]$ via the Weyl--Heisenberg group. This construction has automorphism group of order 144 and is built from the eigenbases of the Weyl--Heisenberg operators $X$ and $Z$ acting on the minimal SI-C set. Crucially, the coordinate entries are drawn from $\{0, 1, \omega, \bar\omega\}$---precisely the Eisenstein alphabet $\{0, \pm 1, \pm\omega, \pm\bar\omega\}$ (after accounting for ray equivalence under scalar multiplication). Cabello's simplest KS set is the smallest member of our Eisenstein island.

This identification provides a concrete bridge between the algebraic island classification and the operational significance of KS sets.

\subsubsection{From KS sets to Bell scenarios: the role of context count}\label{sec:ks-bell}

The operational link between KS sets and nonlocal games is made precise by Trandafir and Cabello~\cite{TrandafirCabello2025optimal}, who give an algorithm to convert any KS set into a bipartite perfect quantum strategy (BPQS) with the minimum number of measurement settings. In such a conversion, the \emph{number of bases} (contexts) in the KS set directly determines the number of inputs for each party in the resulting Bell scenario: fewer bases yield a simpler nonlocal game with fewer measurement choices.

For the Eisenstein island, Cabello~\cite{Cabello2025simplest} shows that the 33-vector, 14-basis KS set yields a $5 \times 9$ qutrit--qutrit BPQS (Alice has 5 inputs, Bob has 9), improving on the $7 \times 9$ strategy derived from the Peres 33-vector set (16 bases).

We independently computed B-KS input counts for all six islands using a SAT-based encoding. Following Definition~2 of Cabello~\cite{Cabello2024bipartite}, a pair $(S_A, S_B)$ of basis subsets is \emph{bipartite KS-uncolorable} (B-KS) if there is no assignment $f: (\mathcal{V}_A \cup \mathcal{V}_B) \times (S_A \cup S_B) \to \{0,1\}$ satisfying:
\begin{itemize}
    \item[(I\,$'$)] \emph{Cross-party exclusion}: $f(u,b) + f(u',b') \leq 1$ for each $b \in S_A$, $b' \in S_B$, $u \in b$, $u' \in b'$ with $u \perp u'$;
    \item[(II\,$'$)] \emph{Completeness}: $\sum_{u \in b} f(u,b) = 1$ for each $b \in S_A \cup S_B$,
\end{itemize}
where $\mathcal{V}_A$ ($\mathcal{V}_B$) is the set of all vectors appearing in $S_A$ ($S_B$). Note that $f$ is contextual: the same vector may receive different values in different bases. Crucially, orthogonality constraints are imposed \emph{only between} $S_A$ and $S_B$ bases, not within a single party's basis set. The basis sets $S_A$ and $S_B$ are allowed to overlap. The optimal BPQS minimizes the product $|S_A| \times |S_B|$ (the total number of input pairs in the Bell scenario). We encode conditions (I$'$) and (II$'$) as a SAT instance and search for the $(S_A, S_B)$ pair minimizing this product, following the approach of Trandafir and Cabello~\cite{TrandafirCabello2025optimal}. For islands with $|\mathcal{B}| \leq 16$ (Eisenstein, Peres, $\Z[\sqrt{-2}]$), we perform exhaustive search over all basis subset pairs; for larger islands (CK-31, Heegner-7, Golden), we use randomized greedy search (500 trials), yielding upper bounds on the optimal product.

\begin{table}[ht]
\centering
\caption{B-KS input counts for all six islands. $|S_A| \times |S_B|$ is the minimum product found over all basis subset pairs that are bipartite KS-uncolorable. Results for the Eisenstein ($5 \times 9$) and Peres ($7 \times 9$) islands independently verify the values reported in Cabello~\cite{Cabello2025simplest} (Table~I therein, obtained using the method of~\cite{TrandafirCabello2025optimal}); CK-31 ($8 \times 9$) matches the same table. The last three are new. ``Exact'' indicates exhaustive search over all $\binom{|\mathcal{B}|}{k}$ subset pairs; ``best found'' indicates randomized greedy search (500 trials), yielding upper bounds on the true optimum. The $\Z[\sqrt{-2}]$ result matches Peres, extending the graph isomorphism of Proposition~\ref{prop:isomorphism} to B-KS structure.}
\label{tab:bpqs}
\begin{tabular}{lcccccl}
\toprule
Island & rays & $|\mathcal{B}|$ & $|S_A|$ & $|S_B|$ & Product & Status \\
\midrule
Eisenstein & 33 & 14 & 5 & 9 & \textbf{45} & Exact; verified \\
Peres & 33 & 16 & 7 & 9 & 63 & Exact; verified \\
$\Z[\sqrt{-2}]$ & 33 & 16 & 7 & 9 & 63 & Exact; new ($=$ Peres) \\
Integer (CK-31) & 31 & 17 & 8 & 9 & $\leq$72 & Best found; matches~\cite{Cabello2025simplest} \\
Heegner-7 & 43 & 23 & 9 & 12 & $\leq$108 & Best found; new \\
Golden & 52 & 25 & 12 & 13 & $\leq$156 & Best found; new \\
\bottomrule
\end{tabular}
\end{table}

\noindent Two empirical patterns emerge from Table~\ref{tab:bpqs}. First, $|S_A| + |S_B| \geq |\mathcal{B}|$ in every case examined: the best-found Alice and Bob basis sets nearly \emph{partition} the full basis set (equality holds for Eisenstein, Peres, CK-31, and Golden; Heegner-7 achieves $21/23$). This pattern suggests---but does not prove---that B-KS uncolorability requires the two parties to collectively engage essentially all measurement contexts. Second, the BPQS cost $|S_A| \times |S_B|$ scales roughly as $|\mathcal{B}|^2/4$, with the ratio product$/|\mathcal{B}|^2$ lying between 0.20 (Heegner-7) and 0.25 (Golden) for all six islands. That the Eisenstein island (product$/|\mathcal{B}|^2 = 0.23$, exact) is more efficient than the $\approx 0.25$ achieved by the real-field islands (Peres, CK-31, Golden) may reflect the richer symmetry of $\Z[\omega]$.

The island classification thus gives direct operational meaning to the algebraic substrate. CK-31 achieves the minimum ray count (31) but at the cost of having the most bases among the 33-class islands (17 vs.\ 14). The Eisenstein island sacrifices two rays to gain a sparser basis structure, yielding the simplest Bell scenario. This trade-off---\emph{ray economy vs.\ context economy}---is invisible without the island framework.

The trade-off extends further when CSW invariants (Section~\ref{sec:csw}) and critical bases analysis (Section~\ref{sec:critical-bases}) are included. No single island dominates on all axes: CK-31 is the tightest proof ($\eta = 1.000$) with a nontrivial CSW sandwich, but its 17 bases make it operationally costly. The Eisenstein island has the simplest BPQS and near-maximal tightness ($\eta = 0.929$), but only a marginal CSW advantage. The Heegner-7 island has the richest CSW structure but the largest BPQS cost. This suggests that the choice of algebraic substrate has genuine consequences for quantum information tasks, and that no single notion of ``simplest'' or ``best'' KS set is appropriate---it depends on the operational criterion.

The fact that Cabello's ``simplest'' construction lands precisely on one of our independently discovered islands---and not between them---is strong evidence that the island classification captures genuine algebraic structure rather than an artifact of our search methodology.

\subsection{Graph-theoretic invariants and Bell inequalities}\label{sec:csw}

If the algebraic substrate controls KS-uncolorability, it should also control the \emph{strength} of the contextual advantage. The Cabello--Severini--Winter (CSW) framework~\cite{CabelloSeveriniWinter2014} provides a quantitative test, mapping the orthogonality graph $G$ of a set of rays to a noncontextuality inequality via three invariants:
\[
\alpha(G) \leq \vartheta(G) \leq \alpha^*(G),
\]
where $\alpha(G)$ is the independence number (classical bound), $\vartheta(G)$ is the Lov\'asz theta number (quantum bound, computed by semidefinite programming), and $\alpha^*(G)$ is the fractional packing number (nonsignaling bound, computed by linear programming). A quantum advantage exists when $\vartheta(G) > \alpha(G)$.

We compute these invariants for two representations of each island: the SAT-minimized KS set (minimum cardinality) and the full ray pool (all rays generated by the alphabet before minimization). This comparison reveals a structural mechanism by which cardinality minimization destroys CSW contextual~advantage.

\begin{table}[ht]
\centering
\caption{CSW graph invariants for minimized KS sets. $\alpha$ = independence number (classical), $\vartheta$ = Lov\'asz theta (quantum, SDP via SCS with tolerance $10^{-9}$), $\alpha^*$ = fractional packing (nonsignaling, LP via HiGHS, exact). The ratio $\vartheta/\alpha$ measures the CSW contextual advantage. Equalities $\vartheta = \alpha$ are reported when $|\vartheta - \alpha| < 10^{-6}$.}
\label{tab:csw-min}
\begin{tabular}{lccccc}
\toprule
Island (minimized) & $n$ & $\alpha(G)$ & $\vartheta(G)$ & $\alpha^*(G)$ & $\vartheta/\alpha$ \\
\midrule
Integer (CK-31) & 31 & 11 & 11.71 & 12.00 & \textbf{1.065} \\
Eisenstein & 33 & 12 & 12.05 & 13.00 & 1.004 \\
Peres & 33 & 12 & 12.00 & 12.00 & 1.000 \\
$\Z[\sqrt{-2}]$ & 33 & 12 & 12.00 & 12.00 & 1.000 \\
Heegner-7 & 43 & 16 & 16.00 & 16.00 & 1.000 \\
\bottomrule
\end{tabular}
\end{table}

\begin{table}[ht]
\centering
\caption{CSW graph invariants for full (non-minimized) ray pools. $\vartheta$ computed via SCS (tolerance $10^{-9}$); $\alpha^*$ via HiGHS (exact). The ``auxiliary'' column counts rays that participate in pairwise orthogonalities but not in any complete basis (Definition~\ref{def:auxiliary}).}
\label{tab:csw-pool}
\begin{tabular}{lccccccc}
\toprule
Island (full pool) & $n$ & bases & auxiliary & $\alpha(G)$ & $\vartheta(G)$ & $\alpha^*(G)$ & $\vartheta/\alpha$ \\
\midrule
Integer & 49 & 26 & 0 (0\%) & 17 & 17.70 & 17.00 & 1.041 \\
Eisenstein & 57 & 22 & 0 (0\%) & 18 & 19.34 & 21.00 & 1.074 \\
Heegner-7 & 145 & 42 & 76 (52\%) & 50 & 55.89 & 66.00 & \textbf{1.118} \\
Peres & 49 & 16 & 16 (33\%) & 23 & 23.00 & 23.00 & 1.000 \\
$\Z[\sqrt{-2}]$ & 49 & 16 & 16 (33\%) & 23 & 23.00 & 23.00 & 1.000 \\
\bottomrule
\end{tabular}
\end{table}

\subsubsection{Auxiliary rays and the CSW advantage}

The comparison between Tables~\ref{tab:csw-min} and~\ref{tab:csw-pool} reveals a striking pattern: the Heegner-7 island's full pool has the strongest contextual advantage of any island ($\vartheta/\alpha = 1.118$), yet its minimized 43-vector set is $\vartheta$-perfect ($\vartheta = \alpha$). To explain this, we introduce a structural distinction among rays in an algebraic pool.

\begin{definition}\label{def:auxiliary}
Let $\mathcal{P}$ be a ray pool with orthogonality graph $G = (V, E)$ and set of bases (maximal orthogonal triples) $\mathcal{B}$. A ray $v \in V$ is \textbf{basis-participating} if $v$ belongs to at least one basis $B \in \mathcal{B}$, and \textbf{auxiliary} otherwise. That is, an auxiliary ray has at least one orthogonal neighbor in $G$ but does not belong to any complete orthonormal triple.
\end{definition}

Auxiliary rays contribute edges to the orthogonality graph without contributing to the hypergraph structure that determines KS-uncolorability. Since KS-uncolorability is a property of the basis hypergraph (no consistent $\{0,1\}$ coloring assigning exactly one ``1'' per basis), auxiliary rays are invisible to it. But they are fully visible to the CSW framework, which depends only on the orthogonality \emph{graph}.

\begin{observation}\label{obs:auxiliary}
\textbf{(Auxiliary rays amplify the CSW advantage.)} Among the five algebraic islands, the CSW contextual advantage $\vartheta(G)/\alpha(G)$ of the full ray pool correlates with the proportion of auxiliary rays:
\begin{center}
\begin{tabular}{lccl}
\toprule
Island & Auxiliary fraction & $\vartheta/\alpha$ (pool) & $\vartheta/\alpha$ (minimized) \\
\midrule
Heegner-7 & 52\% & \textbf{1.118} & 1.000 \\
Peres & 33\% & 1.000 & 1.000 \\
$\Z[\sqrt{-2}]$ & 33\% & 1.000 & 1.000 \\
Integer & 0\% & 1.041 & 1.065 \\
Eisenstein & 0\% & 1.074 & 1.004 \\
\bottomrule
\end{tabular}
\end{center}
Auxiliary rays constrain the independence number $\alpha(G)$ (by adding edges that prevent large independent sets) without contributing to the basis structure. The SDP relaxation $\vartheta(G)$ can exploit these additional constraints by distributing weight fractionally across the auxiliary vertices, widening the gap $\vartheta - \alpha$.
\end{observation}

The mechanism operates as follows. In the Heegner-7 pool ($n = 145$), only 69 rays participate in the 42 bases. The remaining 76 auxiliary rays add 522 orthogonal pairs to the graph, including 9 hub vertices of degree 16 (each connected to 11\% of all other vertices). These hubs severely restrict the size of independent sets: only $\alpha/n = 0.345$ of vertices can be simultaneously independent, compared to $0.469$ for the Peres and $\Z[\sqrt{-2}]$ pools (which are $\vartheta$-perfect). The SDP, however, can assign fractional weights summing to $\vartheta/n = 0.386$, yielding a 12\% quantum advantage.

When SAT minimization reduces the Heegner-7 pool to 43 vectors, it removes auxiliary rays first---they are dispensable for KS-uncolorability. But these are precisely the rays that constrained $\alpha$ and enabled the $\vartheta > \alpha$ gap. The minimized set is $\vartheta$-perfect.

\subsubsection{Further patterns}

\begin{enumerate}
    \item \textbf{Full pools consistently outperform minimized sets} for the CSW advantage. The Eisenstein pool jumps from $\vartheta/\alpha = 1.004$ (minimized) to $1.074$ (full pool). The integer pool is the exception: it drops slightly from $1.065$ to $1.041$, suggesting that CK-31 is unusually efficient at preserving CSW structure despite minimization. Notably, the integer and Eisenstein pools have zero auxiliary rays---every ray participates in a basis---so their CSW advantages arise from graph regularity rather than auxiliary edges.

    \item \textbf{The nonsignaling bound is tightly constrained by basis structure.} For the four minimized sets where $\vartheta = \alpha$ (Peres, $\Z[\sqrt{-2}]$, Heegner-7, plus Eisenstein near-equality), we also have $\alpha^* = \alpha$: all three invariants collapse. Only the integer island (CK-31) maintains a nontrivial sandwich $\alpha = 11 < \vartheta = 11.71 < \alpha^* = 12$. In the full pools, the Heegner-7 pool shows the widest separation: $\alpha = 50 < \vartheta = 55.89 < \alpha^* = 66$, with the nonsignaling bound 32\% above the classical bound.

    \item \textbf{KS-uncolorability does not imply CSW violation}: The Peres and $\Z[\sqrt{-2}]$ islands remain $\vartheta$-perfect even in their full pools, despite having 16 auxiliary rays each. The auxiliary ray mechanism is necessary but not sufficient; the specific placement of orthogonalities matters. In these two pools, the orthogonality graph is disconnected: the 16 auxiliary rays (degree 3) form loosely connected components that do not bridge the main structure. For disconnected graphs, $\vartheta(G) = \sum_i \vartheta(G_i)$ and $\alpha(G) = \sum_i \alpha(G_i)$, so $\vartheta$-perfection of each component implies $\vartheta$-perfection of the whole~\cite{Knuth1994}. This contrasts with the Heegner-7 and integer pools, where auxiliary rays create a connected graph with enough edge density to separate $\vartheta$ from $\alpha$.
\end{enumerate}

\begin{remark}
The distinction between basis-participating and auxiliary rays reveals that KS-uncolorability and CSW contextual advantage are driven by different structural features of the same algebraic pool. Uncolorability requires bases arranged so that no valid $\{0,1\}$ coloring exists---a property of the basis \emph{hypergraph}. The CSW advantage requires pairwise orthogonalities arranged so that $\vartheta > \alpha$---a property of the orthogonality \emph{graph}. Auxiliary rays contribute to the latter without affecting the former. The operational lesson is that the contextual strength of an algebraic island, in the CSW sense, is determined by its full ray pool, not its minimal KS subset. Minimizing for cardinality can destroy the CSW advantage entirely.
\end{remark}

\subsubsection{Spectral explanation of the CSW advantage}\label{sec:spectral}

The Hoffman bound provides a spectral explanation for which pools admit a CSW violation. For any graph $G$ with adjacency matrix $A$ having largest eigenvalue $\lambda_{\max}$ and smallest eigenvalue $\lambda_{\min}$, the independence number satisfies~\cite{Haemers1995}
\[
\alpha(G) \leq \frac{n \cdot (-\lambda_{\min})}{\lambda_{\max} - \lambda_{\min}}.
\]
This bound is valid for all graphs (not only regular graphs); for regular graphs $\lambda_{\max}$ equals the common degree, but the inequality holds in general with $\lambda_{\max}$ as the spectral radius. However, for disconnected graphs the bound can be loose, since the spectral radius of each component contributes independently. When $\alpha(G)$ is well below this bound, the graph is ``spectrally tight'': the eigenvalue structure constrains independent sets more than the combinatorial structure alone requires, leaving room for the SDP relaxation $\vartheta(G)$ to exceed $\alpha(G)$.

\begin{table}[ht]
\centering
\caption{Spectral properties of full ray pool orthogonality graphs. $H$ = Hoffman bound on $\alpha$. The ratio $\alpha/H$ measures spectral tightness: values near or below 1 correlate with CSW violations.}
\label{tab:spectral}
\begin{tabular}{lccccccc}
\toprule
Pool & $n$ & $\lambda_{\max}$ & $\lambda_{\min}$ & $H$ & $\alpha$ & $\alpha/H$ & $\vartheta/\alpha$ \\
\midrule
Heegner-7 & 145 & 8.25 & $-4.50$ & 51.2 & 50 & \textbf{0.977} & \textbf{1.118} \\
Eisenstein & 57 & 6.42 & $-3.45$ & 19.9 & 18 & \textbf{0.904} & \textbf{1.074} \\
Integer & 49 & 5.98 & $-3.29$ & 17.4 & 17 & 0.977 & 1.041 \\
Peres & 49 & 5.53 & $-3.68$ & 19.6 & 23 & 1.176 & 1.000 \\
$\Z[\sqrt{-2}]$ & 49 & 5.53 & $-3.68$ & 19.6 & 23 & 1.176 & 1.000 \\
\bottomrule
\end{tabular}
\end{table}

The pattern in Table~\ref{tab:spectral} is clear: the three pools with $\alpha/H < 1$ (spectrally tight) have CSW violations, with the largest violation ($\vartheta/\alpha = 1.118$) occurring for the Heegner-7 pool ($\alpha/H = 0.977$). The two $\vartheta$-perfect pools have $\alpha/H > 1$, meaning $\alpha$ exceeds the Hoffman bound; this occurs because these graphs are disconnected (the 16 auxiliary rays of degree 3 form loosely connected components), invalidating the bound for non-connected graphs.

\subsection{Contextual tightness: critical bases analysis}\label{sec:critical-bases}

The CSW framework measures how much quantum mechanics exceeds classical bounds on the orthogonality \emph{graph}. We now introduce a complementary measure that probes the internal structure of the basis \emph{hypergraph}: how many bases must be removed before a KS set becomes colorable?

\begin{definition}\label{def:essential-basis}
Let $\mathcal{S}$ be a KS set with basis set $\mathcal{B} = \{B_1, \ldots, B_m\}$. A basis $B_i$ is \textbf{essential} if $\mathcal{S}$ with the constraint from $B_i$ removed admits a valid $\{0,1\}$ coloring---i.e., the set is KS-uncolorable only when $B_i$'s ``exactly one green'' rule is enforced. The \textbf{critical number} $\kappa(\mathcal{S})$ is the minimum $k$ such that some $k$-element subset of bases can be removed to make $\mathcal{S}$ colorable. The \textbf{contextual fraction} is $\kappa(\mathcal{S}) / |\mathcal{B}|$.
\end{definition}

Note that removing a basis relaxes the ``exactly one green per triple'' constraint for that context, while all pairwise orthogonality constraints (``at most one green'') remain---orthogonality is a geometric fact, not a contextual choice. A set with $\kappa = 1$ is \emph{fragile}: a single basis removal suffices to restore colorability. A set with large $\kappa$ is \emph{robust}: many independent constraints must be simultaneously relaxed.

We compute these quantities for each island's SAT-minimized KS set and, for comparison, the full ray pools of the Peres and Heegner-7 islands. The \textbf{essentiality ratio} $\eta = (\text{essential bases}) / |\mathcal{B}|$ measures what fraction of bases are individually indispensable.

\begin{table}[ht]
\centering
\caption{Critical bases analysis. $|\mathcal{B}|$ = number of bases, ess = essential bases (whose individual removal breaks uncolorability), $\eta$ = essentiality ratio, $\kappa$ = critical number (minimum bases to remove), CF = contextual fraction $\kappa/|\mathcal{B}|$.}
\label{tab:critical-bases}
\begin{tabular}{lcccccc}
\toprule
Set & rays & $|\mathcal{B}|$ & ess & $\eta$ & $\kappa$ & CF \\
\midrule
Integer (CK-31) & 31 & 17 & \textbf{17} & \textbf{1.000} & 1 & 0.059 \\
Eisenstein min-33 & 33 & 14 & 13 & 0.929 & 1 & 0.071 \\
Peres min-33 & 33 & 16 & 13 & 0.812 & 1 & 0.063 \\
$\Z[\sqrt{-2}]$ min-33 & 33 & 16 & 13 & 0.812 & 1 & 0.063 \\
Heegner-7 min-43 & 43 & 23 & 18 & 0.783 & 1 & 0.044 \\
\midrule
Peres pool & 49 & 16 & 13 & 0.812 & 1 & 0.063 \\
Heegner-7 pool & 145 & 42 & 0 & 0.000 & \textbf{2} & 0.048 \\
\bottomrule
\end{tabular}
\end{table}

Three findings emerge from Table~\ref{tab:critical-bases}:

\begin{enumerate}
    \item \textbf{CK-31 is maximally tight.} Every one of its 17 bases is essential: $\eta = 1.000$. No basis is redundant for the KS proof. This is a structural consequence of CK-31's minimality---at 31 vectors, it is the smallest known KS set in $\R^3$, and it achieves this by having every constraint contribute non-redundantly to the coloring obstruction.

    \item \textbf{Peres and $\Z[\sqrt{-2}]$ are again identical.} Both have 16 bases, 13 essential, 3 removable, and the same essentiality ratio of 0.812. This extends the graph isomorphism of Proposition~\ref{prop:isomorphism} to the hypergraph level: not only do these two islands have isomorphic orthogonality graphs, they have identical critical bases structure.

    \item \textbf{The Heegner-7 full pool is the only set with $\kappa > 1$.} No single basis removal breaks uncolorability---every basis is individually dispensable (all 42 bases are removable, $\eta = 0$). But removing certain \emph{pairs} of bases does break it: we find exactly 24 critical pairs out of $\binom{42}{2} = 861$ possible pairs, a critical pair density of $0.028$. This high redundancy reflects the Heegner-7 pool's large size (145 rays, 42 bases) compared to the other islands.
\end{enumerate}

The critical number $\kappa = 1$ for all minimized sets means that SAT minimization produces sets where at least one basis is a ``bottleneck'' for the proof. The transition to $\kappa = 2$ in the Heegner-7 full pool shows that algebraic completeness provides a qualitatively different kind of robustness---contextual redundancy---that minimization destroys.

\subsection{Supporting evidence}\label{sec:supporting}

Three additional computational results provide context for the island classification.

\begin{remark}[Realizability gap]\label{rem:realizability}
The algebraic islands are rare among abstract combinatorial structures. In a large-scale experiment, we generated 30{,}000 random 3-uniform hypergraphs with controlled interlocking across sizes $k = 22$--$30$ (5{,}000 per size) and tested both KS-uncolorability (via SAT) and realizability in $\R^3$ (via L-BFGS-B optimization with 50 restarts). Of these, 10{,}059 were KS-uncolorable, but \emph{none} were realizable---the best residuals were consistently $\mathcal{L} \gtrsim 1.0$, far from the $< 10^{-8}$ threshold achieved by all known realizable KS configurations. The geometry of $\R^3$---specifically, the forced completions and transitive constraints of cross-product closure---eliminates essentially all abstract solutions. This realizability bottleneck is uniform from $k = 22$ (near the lower bound) to $k = 30$ (near the upper bound).
\end{remark}

\begin{remark}[Perturbation sensitivity]\label{rem:perturbation}
The KS property is not robust under coordinate perturbation. Adding independent Gaussian noise $\mathcal{N}(0, \varepsilon^2)$ to the coordinates of CK-31 and testing orthogonality at tolerance $\tau = 10^{-6}$, we find that $\varepsilon = 0.01$ suffices to destroy the KS property in all 50 tested configurations. This is ultimately expected---KS sets depend on exact algebraic orthogonalities---but quantifies the sharpness of the transition from exact to approximate contextuality, complementing the rigidity result of Trandafir and Cabello~\cite{TrandafirCabello2025}.
\end{remark}

\begin{remark}[The KS theorem as a rounding impossibility]\label{rem:rounding}
The KS coloring problem admits a natural relaxation: for the maximally mixed state $\rho = I/3$, the Born rule assigns each ray probability $1/3$, and the triad constraint $p(v_1) + p(v_2) + p(v_3) = 1$ is trivially satisfied. The KS theorem says this \emph{fractional} valuation cannot be consistently \emph{rounded} to $\{0,1\}$ across all triads---a feasibility gap analogous to integrality gaps in combinatorial optimization~\cite{Vazirani2001}. The gap between the lower bound (24) and CK-31 (31) measures the cost geometric realizability imposes on abstract combinatorial constraints: forced cross-product completions and transitive orthogonality constraints in $\R^3$ eliminate all candidate structures below 31 vectors (among known algebraic pools), despite the abstract SAT bound permitting solutions at 24.
\end{remark}

\begin{observation}[Universal merge saturation]\label{obs:merge}
For each of the six minimal KS sets, we tested \emph{every} non-orthogonal vertex pair merge (identifying two vertices into one that inherits the union of neighborhoods, discarding degenerate triads). In all 3{,}756 merges across the six islands---394 (CK-31), 456 (Peres-33), 450 (Eisenstein-33), 456 ($\Z[\sqrt{-2}]$-33), 796 (Heegner-7-43), 1{,}204 (Golden-52)---the merged $(n{-}1)$-vertex graph remains KS-uncolorable. We conjecture that every minimal KS set in $\R^3$ or $\C^3$ is \emph{merge-saturated}: merging any non-orthogonal pair preserves KS-uncolorability.
\end{observation}

\begin{conjecture}\label{conj:optimal}
The Conway--Kochen set $\CK$ achieves the true minimum for KS sets in $\R^3$: no KS set with 30 or fewer rays exists.
\end{conjecture}

This conjecture is supported by: the rigidity of $\CK$~\cite{TrandafirCabello2025} (uniqueness suggests optimality); the OCUS proof that no $\leq 30$-ray KS subset exists within the full 49-ray integer pool (Remark after Proposition~\ref{prop:31optimal}); the absence of any sub-31 KS set across all six algebraic islands (Observation~\ref{obs:islands}); universal merge saturation (Observation~\ref{obs:merge}); the complete failure of 30{,}000+ random hypergraphs to achieve geometric realizability; and the graph universality result below.

\begin{observation}[Graph universality of $\CK$]\label{obs:universality}
We tested every algebraic construction in our survey that achieves the minimum of 31 rays---the integer alphabet $\{0,\pm 1,\pm 2\}$, the rational alphabet $\{0,\pm 1,\pm 2,\pm\tfrac{1}{2}\}$, the extended Peres alphabet $\{0,\pm 1,\pm\sqrt{2},\pm(1+\sqrt{2}),\pm 2\}$, the mixed norm-2 alphabet $\{0,\pm 1,\pm 2,\pm\sqrt{2},\pm\omega,\pm\bar\omega\}$ combining three algebraic cancellation identities, and finite group orbits under the tetrahedral group $A_4$ and octahedral group $S_4$ acting on seed vectors in $\R^3$---and in every case the resulting minimal 31-ray KS set has \emph{identical} graph invariants: 71 orthogonality edges, 17 bases, degree sequence $(3^4, 4^{14}, 5^8, 6^3, 8^2)$, identical spectra, and identical Weisfeiler--Leman color refinement signatures. Despite arising from algebraically independent constructions, all roads to 31 lead to the same orthogonality graph.
\end{observation}

This universality strengthens Conjecture~\ref{conj:optimal}: $\CK$ is not merely the smallest KS set found in the integer alphabet, but a unique graph-theoretic attractor across all known algebraic and group-theoretic constructions in $\R^3$. No alternative 31-vertex orthogonality graph has been found despite systematic search across cubic fields $\Q(\sqrt[3]{d})$, Gaussian integers, roots of unity, mixed-field alphabets, and finite rotation group orbits.  The 31-vertex graph appears to be the unique minimal realizable KS hypergraph in dimension~3.

\section{Limitations and Open Questions}\label{sec:limitations}

Our search is computational, not exhaustive over all possible number fields.  It is important to understand both its scope and its principled structure.

The relationship between alphabets, cancellation identities, and KS sets is tight: an alphabet is a candidate for producing a KS set \emph{only if} it carries a cancellation identity---an algebraic relation among its elements that makes dot products vanish exactly.  Without such an identity, the alphabet cannot produce exact orthogonalities among triples of rays, and without orthogonal triples there is no KS set.  The search for new KS-supporting alphabets is therefore equivalent to the search for new cancellation identities of sufficiently low norm.  Our survey is systematic within this framework: for each class of number fields, we identify the candidate cancellation identities from the ring structure and test whether they generate KS-uncolorable ray sets.

The coverage is thorough for low-degree extensions but becomes increasingly sparse at higher degree:

\begin{enumerate}
    \item \textbf{Quadratic fields (near-exhaustive)}: we tested real $\Q(\sqrt{d})$ for all non-square $d = 2, \ldots, 30$ with the basic alphabet $\{0, \pm 1, \pm\sqrt{d}\}$, and also the ring-of-integers generator $(1+\sqrt{d})/2$ for $d \equiv 1\pmod{4}$ ($d \in \{5, 13, 17, 21, 29\}$).  The golden ratio case ($d = 5$) demonstrates that the choice of generator within a field matters: $\varphi = (1+\sqrt{5})/2$ produces a KS set while $\sqrt{5}$ does not, because $\varphi$ participates in a cancellation identity ($\varphi^2 = \varphi + 1$) that $\sqrt{5}$ does not.
    \item \textbf{Complex quadratic fields (exhaustive for class number~1)}: we tested all nine imaginary quadratic fields with class number 1 (the Heegner numbers: $d = 1, 2, 3, 7, 11, 19, 43, 67, 163$), as well as $\Z[\sqrt{-d}]$ for $d = 2, 5, 6, 7, 10, 11, 13, 14, 15$ with enriched alphabets and Hermitian completion.  Higher class numbers (where the ring of integers is no longer a UFD) remain untested; different factorization structures could in principle support cancellation identities not available in class-1 rings.
    \item \textbf{Cyclotomic fields (good coverage)}: all roots of unity of order $n \leq 30$ were tested.  The $6 \mid n$ pattern for KS-uncolorability was verified but not proved for all $n$.
    \item \textbf{Cubic and higher extensions (sparse)}: we tested cubic extensions $\Q(\sqrt[3]{d})$ for $d = 2, 3, 4, 5, 6, 7$. The basic alphabets $\{0,\pm 1,\pm\sqrt[3]{d}\}$ are all colorable, but the extended alphabet $\{0,\pm 1,\pm\sqrt[3]{2},\pm\sqrt[3]{4}\}$ of $\Q(\sqrt[3]{2})$ produces an uncolorable set (minimum 60 vectors after cross-product completion), showing that cubic fields can support KS sets at higher cost.  Quartic, quintic, and higher-degree extensions are entirely unexplored; these could harbor cancellation identities qualitatively different from the quadratic norm-2 pattern.  This is the principal gap in our coverage.
    \item \textbf{Exotic cancellation candidates}: several families of multi-element alphabets from higher-degree fields remain untested and could in principle harbor qualitatively new cancellation mechanisms:
    \begin{itemize}
        \item \emph{Cubic minimal-polynomial cancellations.}  The plastic ratio $\rho \approx 1.3247$ (smallest Pisot number, root of $x^3 = x + 1$) provides a three-element alphabet $\{0, \pm 1, \pm\rho, \pm\rho^2\}$ with cancellation identity $\rho \cdot \rho^2 - \rho - 1 = 0$.  This is neither norm-2 nor phase: the cancellation \emph{is} the minimal polynomial.  Unlike $\Q(\sqrt[3]{2})$ where $\sqrt[3]{2}\cdot\sqrt[3]{4} = 2$ reduces to norm-2, the plastic ratio identity is irreducibly cubic.  Other cubic Pisot numbers (the supersilver ratio, root of $x^3 = 2x^2 + 1$; the tribonacci constant, root of $x^3 = x^2 + x + 1$) provide analogous candidates with different product~algebras.
        \item \emph{Class-number $> 1$ imaginary quadratic fields.}  Our survey tested all nine class-1 fields (Heegner numbers) exhaustively, but rings with non-unique factorization could support cross-cancellation mechanisms unavailable in UFDs.  For example, in $\Z[\sqrt{-5}]$ (class number~2), $2 \cdot 3 = (1{+}\sqrt{-5})(1{-}\sqrt{-5})$; such multiple factorization paths create algebraic relations with no class-1~analogue.
        \item \emph{Simultaneous conjugate cancellations.}  Multi-element alphabets containing both $\alpha$ and $\bar\alpha$ as generators, where these participate in \emph{independent} cancellation identities, could produce richer orthogonality structures than single-generator alphabets.  The Gaussian case already exhibits this ($|i|^2 = 1$ and $|1{+}i|^2 = 2$ coexist), but higher-degree fields offer more conjugate~pairs.
        \item \emph{CM elliptic curve endomorphism rings.}  The $j$-invariants of CM elliptic curves are algebraic integers in large number fields whose Hilbert class polynomials provide exotic minimal-polynomial cancellations of arbitrarily high degree.  This is highly speculative but represents the logical limit of the algebraic approach.
    \end{itemize}
    Of these, the cubic Pisot numbers and class-number $> 1$ fields are the most accessible computationally; preliminary scripts for testing these are included in the repository.
    \item \textbf{Completion termination}: the cross-product completion operation (Section~\ref{sec:closure}) terminates empirically for all tested alphabets, but we do not prove termination in general. For alphabets whose coordinate ring is not closed under the cofactor formula's operations, the completed pool could in principle grow without bound. In practice, we impose a maximum iteration limit and verify stabilization.
    \item \textbf{Transcendental coordinates}: our methods require algebraic coordinates (for exact orthogonality testing); KS sets with transcendental coordinates are not excluded.
    \item \textbf{The $6 \mid n$ conjecture}: verified for $n \leq 30$ but not proved for all $n$. A proof would likely follow from the representation theory of cyclic groups acting on $\C^3$.
    \item \textbf{Minimization is heuristic}: our randomized greedy minimization does not certify global optimality. The ``smallest found'' values reported in our tables are upper bounds on the true minimum KS set size within each island's completed ray pool. For the three 33-vector islands (Peres, Eisenstein, $\Z[\sqrt{-2}]$), we verified that no 32-vector KS subset exists within any known minimal 33-set: across 2{,}000 trials per island, we found 1--3 distinct 33-sets each, and tested all possible single-vector removals---all were colorable. Similarly, for the Heegner-7 island, all 258 subsets of size 42 across 6 distinct 43-sets were colorable. These checks establish criticality (no ray can be removed) but not global optimality within the full pool. For the integer island, we go further: the OCUS procedure (Remark after Proposition~\ref{prop:31optimal}) exhaustively proves that no $\leq 30$-ray KS subset exists within the full 49-ray pool, certifying 31 as the exact minimum for the $\{0,\pm 1,\pm 2\}$ alphabet.
    \item \textbf{Lower bound improvement}: our work does not improve the lower bound of 24; it constrains where a sub-31 construction could come from.  However, these constraints are directly applicable to the realizability-based programs of Li, Bright, and Ganesh~\cite{LiBrightGanesh2024} and Kirchweger et al.~\cite{KirchwegarPeitlSzeider2023}, which work from the opposite direction: they enumerate abstract KS hypergraphs and check realizability in $\R^3$ via SMT solvers.  Our results provide a filter pipeline for their search.  The two-element alphabet cancellation classification (Observation~\ref{obs:two-element}) constrains which number fields can serve as constructive candidates; the density trend (simpler identities produce denser triad networks) sets a structural minimum on the basis count of viable candidate hypergraphs; and the discreteness of the cancellation landscape---the minimal alphabet $\{0, \pm 1\}$ generates only 13 rays, jumping to 49 at $\{0, \pm 1, \pm 2\}$ with no intermediate step---constrains the algebraic degree of coordinates that the SMT solver must handle.  Combining our algebraic pruning with their realizability checking could reduce the search space for sub-31 candidates by orders of magnitude.
    \item \textbf{Operational significance}: our connection to Cabello's work (Section~\ref{sec:cabello}) is primarily observational---we identify the Eisenstein island with his simplest KS set, but do not prove that other islands cannot arise as engines of perfect quantum strategies. A systematic study of which Bell scenarios' perfect strategies correspond to which algebraic islands would deepen this connection.
    \item \textbf{Open questions}: Does CK-31 serve as the engine of any perfect quantum strategy? What BPQS does the Heegner-7 island produce? Can the three-way trade-off (CSW advantage, hypergraph tightness, context economy) be formalized as a Pareto frontier across algebraic islands?
\end{enumerate}

\section{Conclusion}\label{sec:conclusion}

The central empirical finding of this paper is that, among all tested coordinate alphabets, KS-uncolorability in dimension~3 requires one of exactly two cancellation mechanisms: \emph{norm-2 cancellation} (an algebraic identity producing the value~2 from products of ring elements, as in $1+1=2$, $(\sqrt{2})^2=2$, or $\alpha\bar\alpha=2$) or \emph{phase cancellation} (a root-of-unity sum vanishing exactly, as in $1+\omega+\omega^2=0$).  Among all quadratic fields, cyclotomic fields, and selected extensions tested, KS sets exist only when the coordinate ring supports one of these two mechanisms; alphabets where the smallest available cancellation involves norm~$\geq 3$ produce orthogonal triples but not KS-uncolorability.  This two-mechanism constraint explains why constructions cluster into six discrete algebraic islands among the fields tested.  The golden ratio island ($|\varphi|^2 \approx 2.618$) does not violate this boundary: cross-product completion introduces norm-2-type cancellations into the effective coordinate set, confirming that the relevant criterion is the cancellation structure of the completed alphabet, not the generator's own norm.

This structural perspective yields both new constructions and new connections:

\begin{enumerate}
    \item \textbf{New KS configurations.} The Heegner-7 ring $\Z[(1+\sqrt{-7})/2]$ produces a genuinely new orthogonality graph (43 vectors, 23 bases), and the golden ratio field $\Q(\varphi)$ yields a KS set (52 vectors) that is invisible to standard alphabet searches and revealed only by cross-product completion. Neither has been previously reported. The complex quadratic $\Z[\sqrt{-2}]$ provides a new algebraic realization of the Peres-type norm-2 structure, graph-isomorphic to the original (Proposition~\ref{prop:isomorphism}).

    \item \textbf{Operational consequences.} The algebraic substrate determines the complexity of the resulting quantum information tasks. Using SAT-based bipartite KS-uncolorability, we verify and extend the BPQS input counts of Trandafir and Cabello (Table~\ref{tab:bpqs}): the Eisenstein island yields the simplest Bell scenario ($5 \times 9$), while CK-31 is costlier ($\leq 8 \times 9$) despite having fewer rays. CSW analysis shows that auxiliary rays amplify contextual advantage, and that no single island dominates across all operational measures---ray economy, context economy, and CSW strength are in tension.

    \item \textbf{Supporting evidence.} The $6 \mid n$ pattern for cyclotomic uncolorability connects to the $\Z[1/6]$ minimality result of Cortez, Morales, and Reyes; moreover, computing their $N(S)$ invariant for each island shows the Eisenstein island achieves exactly $N = 6$, confirming it as the algebraic realization of the minimal ring.  An OCUS (Optimal Constrained Unsatisfiable Subset) procedure exhaustively proves that no $\leq 30$-ray KS subset exists within the full 49-ray integer pool.  All six minimal KS sets are merge-saturated: every non-orthogonal vertex merge preserves KS-uncolorability (3{,}756 merges tested, 100\% preservation; Observation~\ref{obs:merge}).  A trigonometric consistency check across 164 angles finds only known islands, with KS-uncolorability at isolated algebraic points (Remark~\ref{rem:trig}). Most strikingly, the graph universality result (Observation~\ref{obs:universality}) shows that every construction achieving 31 rays---integers, rationals, mixed-field alphabets, and finite rotation group orbits---produces the \emph{same} orthogonality graph, suggesting that $\CK$ is the unique minimal realizable KS hypergraph in dimension~3.
\end{enumerate}

Whether the two-mechanism constraint extends to all number fields, or whether higher-degree extensions harbor additional cancellation mechanisms and hence additional islands, remains the principal open question. The identification of Cabello's ``simplest KS set'' with the Eisenstein island suggests that the algebraic perspective may complement the ongoing search for optimal constructions. More concretely, the constraints established here---the two-element cancellation classification, density monotonicity, and the discrete 13$\to$49 ray jump---are directly applicable to the realizability-based lower-bound program of Li, Bright, and Ganesh~\cite{LiBrightGanesh2024}: algebraic pruning of candidate hypergraphs could substantially narrow the search space that their SMT solver must explore, bringing the constructive upper bound (31) and the realizability lower bound (24) closer together.

There is a pleasing circularity in this enterprise: the KS theorem establishes the contextuality that underlies quantum computational advantage~\cite{BravyiGossetKoenig2018}, yet resolving the minimum KS set problem may itself require quantum resources.  The SAT and SMT instances arising from KS realizability checking are natural candidates for quantum speedup, suggesting that the very quantum effects whose foundations we probe may ultimately provide the computational power needed to complete the~proof.

\subsection*{Reproducibility}

\begin{sloppypar}
All code and data are available at
\url{https://github.com/michaelkernaghan/contextuality}.
The repository includes explicit ray coordinates and triad/basis lists for all islands, including the 43-vector Heegner-7 and 52-vector golden-ratio sets, as well as verifier scripts that reconstruct each KS set from its alphabet and confirm UNSAT via SAT solving. The implementation uses Python~3.11 with PySAT~1.8 (Glucose4), NumPy, SciPy (L-BFGS-B and HiGHS~LP), and CVXPY (SCS~SDP, tolerance~$10^{-9}$).
\end{sloppypar}

The script-to-table correspondence is:
\begin{itemize}\setlength{\itemsep}{0pt}\setlength{\parskip}{0pt}
\item \texttt{ks\_islands.py} $\to$ Tables~\ref{tab:alphabets}--\ref{tab:closure}
\item \texttt{ks\_complex.py} $\to$ Table~\ref{tab:roots}
\item \texttt{ks\_new\_islands.py}, \texttt{ks\_sat.py} $\to$ Table~\ref{tab:complex-quadratic}, 31-optimality check
\item \texttt{ks\_new\_island\_analysis.py} $\to$ Table~\ref{tab:heegner}
\item \texttt{ks\_csw.py}, \texttt{ks\_csw\_extended.py}, \texttt{ks\_graph\_analysis.py} $\to$ Tables~\ref{tab:csw-min}--\ref{tab:csw-pool},~\ref{tab:spectral}
\item \texttt{ks\_critical\_bases.py} $\to$ Table~\ref{tab:critical-bases};\enspace \texttt{ks\_bpqs\_sat.py} $\to$ Table~\ref{tab:bpqs}
\item \texttt{ks\_trig.py} $\to$ Remark~\ref{rem:trig};\enspace \texttt{ks\_graph\_analysis.py} $\to$ Prop.~\ref{prop:isomorphism}
\item \texttt{ks\_group\_isomorphism.py} $\to$ Obs.~\ref{obs:universality};\enspace \texttt{ks\_explore\_new*.py} $\to$ cubic/mixed/orbit explorations
\item \texttt{ks\_maxsat\_optimal.py} $\to$ OCUS proof (Remark after Prop.~\ref{prop:31optimal});\enspace \texttt{ks\_peres\_merge.py} $\to$ Obs.~\ref{obs:merge}
\item \texttt{ks\_literature\_connections.py} $\to$ SI-C closure comparison, $N(S)$ invariant computation
\end{itemize}
All randomized experiments use \texttt{random.seed(42)} for~reproducibility.

\subsection*{Acknowledgments}

The author thanks Ad\'an Cabello for pointing out the connection between the Eisenstein island and his ``simplest KS set'' construction. The author acknowledges the use of Claude (Anthropic) for computational assistance, literature synthesis, and manuscript preparation.

\begin{thebibliography}{99}

\bibitem{KochenSpecker1967}
S.~Kochen and E.~P.~Specker,
``The problem of hidden variables in quantum mechanics,''
\textit{J. Math. Mech.} \textbf{17}, 59--87 (1967).

\bibitem{Peres1991}
A.~Peres,
``Two simple proofs of the Kochen--Specker theorem,''
\textit{J. Phys. A: Math. Gen.} \textbf{24}, L175--L178 (1991).

\bibitem{Peres1993}
A.~Peres,
\textit{Quantum Theory: Concepts and Methods}
(Kluwer Academic Publishers, 1993).

\bibitem{Kernaghan1994}
M.~Kernaghan,
``Bell--Kochen--Specker theorem for 20 vectors,''
\textit{J. Phys. A: Math. Gen.} \textbf{27}, L829--L830 (1994).

\bibitem{KernaghanPeres1995}
M.~Kernaghan and A.~Peres,
``Kochen--Specker theorem for eight-dimensional space,''
\textit{Phys. Lett. A} \textbf{198}, 1--5 (1995).

\bibitem{CabelloEstebaranzGarcia1996}
A.~Cabello, J.~M.~Estebaranz, and G.~Garc\'{\i}a-Alcaine,
``Bell--Kochen--Specker theorem: A proof with 18 vectors,''
\textit{Phys. Lett. A} \textbf{212}, 183--187 (1996).

\bibitem{Pavicic2004}
M.~Pavicic, J.-P.~Merlet, B.~McKay, and N.~D.~Megill,
``Kochen--Specker vectors,''
\textit{J. Phys. A: Math. Gen.} \textbf{38}, 1577--1592 (2005).
arXiv:quant-ph/0409014.

\bibitem{Pavicic2019}
M.~Pavicic, N.~D.~Megill, B.~C.~Waegell, and P.~K.~Aravind,
``Automated generation of Kochen--Specker sets,''
\textit{Sci. Rep.} \textbf{9}, 6765 (2019).
arXiv:1905.01009.

\bibitem{CabelloSeveriniWinter2014}
A.~Cabello, S.~Severini, and A.~Winter,
``Graph-theoretic approach to quantum correlations,''
\textit{Phys. Rev. Lett.} \textbf{112}, 040401 (2014).

\bibitem{UijlenWesterbaan2016}
S.~Uijlen and B.~Westerbaan,
``A Kochen--Specker system has at least 22 vectors,''
\textit{New Generation Computing} \textbf{34}, 3--23 (2016).

\bibitem{AudemardSimon2018}
G.~Audemard and L.~Simon,
``On the Glucose SAT solver,''
\textit{Int. J. Artif. Intell. Tools} \textbf{27}, 1840001 (2018).

\bibitem{KunjwalSpekkens2018}
R.~Kunjwal and R.~W.~Spekkens,
``Beyond the Cabello--Severini--Winter framework: Making sense of contextuality without sharpness of measurements,''
\textit{Phys. Rev. A} \textbf{97}, 052110 (2018).

\bibitem{CortezMoralesReyes2022}
V.~Cortez, R.~Morales, and E.~J.~Reyes,
``Minimal ring extensions of the integers exhibiting Kochen--Specker contextuality,''
arXiv:2211.13216 (2022).

\bibitem{LiBrightGanesh2024}
J.~Li, C.~Bright, and V.~Ganesh,
``Leveraging SAT solvers for the minimum Kochen--Specker problem,''
in \textit{Proceedings of the 33rd International Joint Conference on Artificial Intelligence (IJCAI-24)}, pp.~1944--1952 (2024).
arXiv:2306.13319.

\bibitem{Cabello2024bipartite}
A.~Cabello,
``Simplest bipartite perfect quantum strategies,''
\textit{Phys.\ Rev.\ Lett.} \textbf{134}, 010201 (2025).
arXiv:2311.17735.

\bibitem{Cabello2024equivalences}
Y.~Liu, H.~Y.~Chung, E.~Z.~Cruzeiro, J.~R.~Gonzales-Ureta, R.~Ramanathan, and A.~Cabello,
``Equivalence between face nonsignaling correlations, full nonlocality, all-versus-nothing proofs, and pseudotelepathy,''
\textit{Phys.\ Rev.\ Research} \textbf{6}, L042035 (2024).
arXiv:2310.10600.

\bibitem{Cabello2025simplest}
A.~Cabello,
``Simplest Kochen--Specker set,''
\textit{Phys.\ Rev.\ Lett.} (2025).
arXiv:2508.07335.

\bibitem{TrandafirCabello2025}
S.~Trandafir and A.~Cabello,
``Two fundamental solutions to the rigid Kochen--Specker set problem and the solution to the minimal Kochen--Specker set problem under one assumption,''
\textit{Phys.\ Rev.\ A} \textbf{111}, 052204 (2025).
arXiv:2501.11640.

\bibitem{TrandafirCabello2025optimal}
S.~Trandafir and A.~Cabello,
``Optimal conversion of Kochen--Specker sets into bipartite perfect quantum strategies,''
\textit{Phys.\ Rev.\ A} \textbf{111}, 022408 (2025).
arXiv:2410.17470.

\bibitem{AbramskyBarbosaMansfield2017}
S.~Abramsky, R.~S.~Barbosa, and S.~Mansfield,
``The contextual fraction as a measure of contextuality,''
\textit{Phys.\ Rev.\ Lett.} \textbf{119}, 050504 (2017).
arXiv:1705.07918.

\bibitem{Vazirani2001}
V.~V.~Vazirani,
\textit{Approximation Algorithms}
(Springer, 2001).

\bibitem{Gleason1957}
A.~M.~Gleason,
``Measures on the closed subspaces of a Hilbert space,''
\textit{J. Math. Mech.} \textbf{6}, 885--893 (1957).

\bibitem{Haemers1995}
W.~H.~Haemers,
``Interlacing eigenvalues and graphs,''
\textit{Linear Algebra Appl.} \textbf{226--228}, 593--616 (1995).

\bibitem{Knuth1994}
D.~E.~Knuth,
``The sandwich theorem,''
\textit{Electron.\ J.\ Combin.} \textbf{1}, Article A1 (1994).

\bibitem{GouldAravind2010}
P.~K.~Gould and P.~K.~Aravind,
``Isomorphism between the Peres and Penrose proofs of the BKS theorem in three dimensions,''
\textit{Found.\ Phys.} \textbf{40}, 1096--1101 (2010).

\bibitem{KirchwegarPeitlSzeider2023}
M.~Kirchweger, A.~Peitl, and S.~Szeider,
``Co-certificate learning with SAT modulo symmetries,''
in \textit{Proceedings of AAAI-23}, 4119--4126 (2023).

\bibitem{Budroni2022}
C.~Budroni, A.~Cabello, O.~G\"uhne, M.~Kleinmann, and J.-\AA.~Larsson,
``Kochen--Specker contextuality,''
\textit{Rev.\ Mod.\ Phys.} \textbf{94}, 045007 (2022).

\end{thebibliography}

\appendix

\section{Representative minimal KS sets for new islands}\label{app:rays}

For reproducibility, we describe the minimal KS sets found for the two genuinely new orthogonality graph types.

\textbf{Heegner-7 island} (43 vectors, 23 bases). Let $\alpha = (1+\sqrt{-7})/2$. All ray coordinates are drawn from $\{0, \pm 1, \pm\alpha, \pm\bar\alpha\}$, with cross-product completion adding elements of $\Z[\alpha]$ (integer linear combinations of $1$ and $\alpha$). The 145-ray pool is generated by enumerating all projectively distinct nonzero vectors $(a, b, c) \in \{0, \pm 1, \pm\alpha, \pm\bar\alpha\}^3$ under Hermitian orthogonality, then closing under cross products. The 43-ray minimal set is obtained by SAT-based greedy minimization (1000 trials); all single-ray removals from this set restore colorability. The 23 bases (orthogonal triples) and the degree distribution $\{4{:}16,\; 5{:}20,\; 6{:}4,\; 8{:}2,\; 10{:}1\}$ distinguish this hypergraph from all other known 3D KS configurations.

\textbf{Golden island} (52 vectors, 25 bases). Let $\varphi = (1+\sqrt{5})/2$. The raw alphabet $\{0, \pm 1, \pm\varphi\}$ generates 49 rays that are colorable; cross-product completion adds rays with coordinates in $\Z[\varphi] = \{a + b\varphi : a, b \in \Z\}$, expanding the pool to 205 rays. The completed pool is KS-uncolorable, with the smallest subset found containing 52 vectors in 25 bases. Completion-generated coordinates include $\varphi^2 = \varphi + 1$ and $2\varphi - 1 = \sqrt{5}$, both elements of $\Z[\varphi]$.

Complete ray coordinates (in exact algebraic form), basis hyperedge lists, and SAT verification scripts for all six islands are provided in the accompanying code repository.

\end{document}
